% !TEX program = tectonic
\documentclass[preprint,12pt]{elsarticle}
\biboptions{numbers}

\usepackage[T1]{fontenc}
\usepackage{microtype}
\usepackage{amsmath,amssymb,amsthm}
\usepackage{newtxtext}
\usepackage{newtxmath}
\renewcommand{\sfdefault}{\rmdefault}
\renewcommand{\ttdefault}{\rmdefault}
\usepackage{booktabs}
\usepackage{tabularx}
\usepackage{array}
\usepackage{xcolor}
\usepackage{graphicx}
\usepackage{float}
\usepackage{placeins}
\usepackage{caption}
\usepackage{subcaption}
\captionsetup[figure]{name=Figure,labelsep=colon}
\usepackage{tikz}
\usetikzlibrary{arrows.meta,positioning,fit,calc,shapes.geometric}
\usepackage[ruled,linesnumbered,noend]{algorithm2e}
\SetAlgorithmName{Algorithm}{Algorithm}{List of Algorithms}
\SetAlgoCaptionSeparator{:}
\SetKwComment{Comment}{$\triangleright$~}{}
\usepackage[colorlinks,citecolor=blue!70!black,linkcolor=blue!70!black,urlcolor=blue!50!black]{hyperref}
\urlstyle{same}
\usepackage[capitalise,noabbrev]{cleveref}
\crefname{figure}{Figure}{Figures}
\Crefname{figure}{Figure}{Figures}

\newcolumntype{Y}{>{\raggedright\arraybackslash}X}
\newcolumntype{L}[1]{>{\raggedright\arraybackslash}p{#1}}
\setlength{\emergencystretch}{3em}
\hfuzz=20pt
\vfuzz=20pt
\hbadness=10000
\vbadness=10000
\raggedbottom
\setlength{\textfloatsep}{7pt plus 1pt minus 2pt}
\setlength{\floatsep}{6pt plus 1pt minus 1pt}
\setlength{\intextsep}{7pt plus 1pt minus 1pt}
\setlength{\bibsep}{1pt plus 0.2ex}
\setcounter{topnumber}{5}
\setcounter{bottomnumber}{4}
\setcounter{totalnumber}{8}
\renewcommand{\topfraction}{0.95}
\renewcommand{\bottomfraction}{0.85}
\renewcommand{\textfraction}{0.05}
\renewcommand{\floatpagefraction}{0.75}
\makeatletter
\setlength{\@fptop}{0pt}
\setlength{\@fpsep}{8pt plus 1fil}
\setlength{\@fpbot}{0pt plus 1fil}
\makeatother

\newtheorem{theorem}{Theorem}[section]
\newtheorem{lemma}[theorem]{Lemma}
\newtheorem{claim}[theorem]{Claim}
\theoremstyle{definition}
\newtheorem{definition}[theorem]{Definition}
\newtheorem{remark}[theorem]{Remark}

\newcommand{\tech}[1]{\textnormal{#1}}
\newcommand{\domain}[1]{\textnormal{``#1''}}
\renewcommand{\mathsf}[1]{\mathrm{#1}}
\newcommand{\Aura}{\textsc{Aura}}
\newcommand{\MLS}{\tech{MLS}}
\newcommand{\TreeKEM}{\tech{TreeKEM}}
\newcommand{\Shield}{\tech{Shield}}
\newcommand{\HKDF}{\tech{HKDF}}
\newcommand{\HMAC}{\tech{HMAC}}
\newcommand{\SHA}{\tech{SHA-256}}
\newcommand{\AEAD}{\tech{AES-256-GCM-SIV}}
\newcommand{\BLAKE}{\tech{BLAKE2b}}
\newcommand{\ProVerif}{\tech{ProVerif}}
\newcommand{\Tamarin}{\tech{Tamarin}}
\newcommand{\EdDSA}{\tech{Ed25519}}
\renewcommand{\DH}{\tech{X25519}}
\newcommand{\KEM}{\tech{ML-KEM-768}}
\newcommand{\Enc}{\tech{Enc}}
\newcommand{\Dec}{\tech{Dec}}
\newcommand{\KeyGen}{\tech{KeyGen}}
\newcommand{\Encaps}{\tech{Encaps}}
\newcommand{\Decaps}{\tech{Decaps}}
\newcommand{\concat}{\,\|\,}
\newcommand{\gid}{\mathit{gid}}
\newcommand{\epoch}{\mathit{epoch}}
\newcommand{\gch}{\mathit{gch}}
\newcommand{\initsec}{\mathit{initSecret}}
\newcommand{\commitsec}{\mathit{commitSecret}}
\newcommand{\epochsec}{\mathit{epochSecret}}
\newcommand{\policy}{\mathit{policy}}
\newcommand{\treehash}{\mathit{treeHash}}

\date{}

\makeatletter
\def\ps@pprintTitle{%
  \let\@oddhead\@empty
  \let\@evenhead\@empty
  \let\@oddfoot\@empty
  \let\@evenfoot\@empty}
\makeatother

\begin{document}

\begin{frontmatter}

\title{Aura Group: Policy-Bound Hybrid TreeKEM for Post-Quantum Group Messaging}

\author{Oleksandr Melnychenko}

\address{Khmelnytskyi National University, Khmelnytskyi, Ukraine}

\journal{Journal of Information Security and Applications}

\begin{abstract}
\Aura{} Group is the group-messaging protocol for \Aura{} Messenger. It
targets dynamic post-quantum confidential groups in which membership changes,
application-message traffic, relay-blind delivery, and local state persistence
belong to the same cryptographic state machine. The design is inspired by
\MLS{} and \TreeKEM{}, but it is not wire-compatible with MLS.

Accepted group epochs advance through \TreeKEM{}-style Commit messages whose
update paths carry both \DH{} and \KEM{} material. Path secrets are encrypted
to co-path resolutions with a hybrid \DH{}+\KEM{} wrapping key, and every
accepted Commit derives a group-epoch secret from the previous init secret,
the fresh commit secret, and a context hash over the group identifier, epoch
number, tree hash, and serialized security policy. Policy is therefore
cryptographic input: disagreement about \Shield{}, external joins, franking,
sender-key lifetime, or skipped-key storage changes the accepted epoch state.
Within an accepted epoch, per-member sender chains provide low-cost
application messages while sender data, payloads, replay state, sealed
content, expiring content, and franking commitments remain bound to the same
epoch and policy.

The security account targets membership-change forward secrecy,
post-compromise recovery at group-epoch boundaries, policy consensus, replay
rejection, and accountable franking. Evidence is deliberately scoped: a
\ProVerif{} model discharges six queries for one accepted hybrid group epoch
and one application message, while the reference implementation validates
parsing, persistence, replay, external-join rejection, \Shield{} enforcement,
franking, sealed state, and benchmarked cost centers. Full asynchronous group
security, multi-device state, delivery fairness, post-quantum signatures, and
arbitrary proposal interleavings remain separate proof and deployment
obligations.
\end{abstract}

\begin{keyword}
secure group messaging \sep TreeKEM \sep ML-KEM \sep message franking
\end{keyword}

\end{frontmatter}

\section{Introduction}

End-to-end encrypted group messaging protects a moving collective state, not
only a sequence of pairwise messages. The central object is an accepted
\emph{group epoch}: a membership interval in which the roster, ratchet-tree
hash, security policy, and application-secret schedule are fixed for all
honest members. When membership changes, the protocol must move the group into
a new epoch, deliver fresh secrets only to the new roster, prevent removed
members from following future traffic, and keep ordinary application messages
efficient between such state transitions.

\Aura{} Group is the group-messaging protocol for \Aura{} Messenger. The
target setting includes ordinary private conversations, but also long-lived
operational groups in which messages can carry inspection evidence, sensor
alerts, incident decisions, or other records whose confidentiality matters
after the immediate exchange has ended. In that setting, payload encryption
alone is not the full security question. A relay must not become an authority
over membership or content; a removed member must not learn later traffic; a
re-added member must not recover messages from the absence interval; and
stored ciphertext should remain protected against future post-quantum
adversaries. This is the common threat shape behind coordinated operational
communication: membership changes, evidence reporting, delayed review, and
untrusted relay delivery occur in the same channel. Prior work on UAV
trajectory adaptation, multimodal wind-turbine inspection, and
UAV/SCADA-linked fire detection gives examples of domains with this
coordination pattern~\cite{svystun2025dytam,svystun2024thermal,lysyi2025fire}.
Those applications are not protocol assumptions. They make concrete the
broader requirement: membership, sender authentication, policy consensus, and
post-quantum confidentiality must be treated as one cryptographic surface.

\MLS{} and \TreeKEM{} provide the closest standard vocabulary for this surface.
An MLS-style Commit changes the authenticated group state, refreshes a path in
the ratchet tree, distributes path secrets to co-path resolutions, and installs
the next epoch after client-side verification~\cite{rfc9420,treekem,treesync}.
The post-quantum migration problem is sharper in this setting than in a
pairwise handshake. A group update that only refreshes classical DH material
leaves the epoch transition dependent on assumptions that may not protect
archived ciphertext against a future cryptographically relevant quantum
computer. The terminology and parameter choices follow the NIST post-quantum
standardization line and the ML-KEM specification used by current hybrid
deployments~\cite{nist-pqc,fips203}.

\Aura{} Group adopts the TreeKEM state-transition discipline but defines a
compact hybrid protocol profile rather than an MLS wire-compatible extension.
The technical question is specific: what changes when ML-KEM material is
carried in TreeKEM update-path wrapping, and when serialized application
security policy is included in the epoch context, key schedule, and
confirmation input? Each hybrid update path carries both \DH{} and \KEM{}
material, and the accepted group context binds the group identifier, epoch
number, ratchet-tree hash, and serialized security policy. The post-quantum
contribution is therefore part of the membership-changing tree update itself,
not only an admission mechanism or a group-creation extension. Policy is also
made cryptographic: a disagreement about \Shield{}, external joins, franking,
sender-key lifetime, or skipped-key storage changes the derived epoch state
instead of remaining an application-level preference.

This design point is intentionally narrower than a new general group-security
framework. The novelty is the protocol interface: hybrid path wrapping,
policy-bound epoch derivation, efficient sender-chain traffic, and accountable
message semantics are placed behind the same local acceptance boundary. That
boundary is important because each mechanism alone is familiar; the security
question is whether the accepted tree, accepted policy, sender-chain state, and
message semantics move together without giving the relay authority over any of
them.

The security account has a defined proof boundary. The specification and
mechanized model focus on one accepted group epoch and on the transition
mechanisms that lead into it. The accompanying implementation evidence checks
parsing, persistence, replay handling, state rejection, and message semantics.
A full concurrent-group theorem would require a separate tree-concurrency
model covering all interleavings of Add, Remove, Update, PSK, and ReInit
proposals.

\medskip
\noindent\textbf{Objective.} This paper specifies and evaluates a dynamic
group-messaging protocol that combines hybrid TreeKEM update paths,
policy-bound epoch derivation, and efficient sender-chain traffic while keeping
the proof, symbolic-model, and implementation-validation boundaries explicit
enough for external review.

\medskip
\noindent\textbf{Contributions and novelty.}
\Aura{} Group contributes a policy-bound hybrid \TreeKEM{} profile in which
post-quantum path refresh, policy consensus, sender-key economy, and
epoch-bound message-layer hooks are specified as one group state machine. The
contributions are as follows.

\begin{enumerate}
  \item \textbf{G1: hybrid TreeKEM update paths.} Every epoch transition can
  refresh path material with both \DH{} and \KEM{} contributions. The new path
  secrets are delivered to co-path resolutions through hybrid wrapping keys,
  so post-quantum material participates in ordinary membership-changing
  updates.

  \item \textbf{G2: policy-bound group context.} The serialized security policy is
  hashed into the group context used for key derivation and confirmation.
  Silent disagreement about \Shield{}, external joins, franking, sender-key
  lifetime, or skipped-key limits therefore changes the epoch keys rather than
  producing an apparently valid downgrade.

  \item \textbf{G3: cryptographic Shield profile.} \Shield{} is defined as a
  protocol profile rather than an advisory product mode. It selects
  profile-specific subkey labels, mandatory franking, external-join blocking,
  bounded skipped-key storage, and bounded sender-key lifetime as inputs to
  the accepted epoch state.

  \item \textbf{G4: sender-chain traffic inside an accepted group epoch.}
  Per-member sender chains provide low-cost application messages between
  TreeKEM Commit messages, while encrypted sender data, replay windows, and
  epoch binding prevent duplicate or cross-interval acceptance.

  \item \textbf{G5: epoch-bound message-layer hooks.} Franking, sealed content,
  and disappearing-message checks are bound to the accepted epoch and policy.
  They are not independent security claims; they specify where higher-level
  message semantics are checked after cryptographic decryption. The relay
  observes only structural metadata and opaque commitments.

  \item \textbf{G6: scoped mechanized evidence.} A dedicated \ProVerif{} model
  checks six obligations for one accepted group epoch: hybrid path secrecy,
  previous-init contribution secrecy, application-message secrecy, Commit
  correspondence, message correspondence, and franking correspondence. The
  scope is stated explicitly, so symbolic evidence is not confused with a full
  proof for arbitrary concurrent group schedules.
\end{enumerate}

\begin{table}[H]
\centering
\caption{Contribution-to-evidence map for the \Aura{} Group scope.}
\label{tab:group-contribution-map}
\footnotesize
\setlength{\tabcolsep}{2pt}
\begin{tabularx}{\textwidth}{@{}L{0.08\textwidth}L{0.30\textwidth}Y@{}}
\toprule
ID & Group claim & Main evidence in this paper \\
\midrule
G1 & Hybrid update paths carry post-quantum material into ordinary group-epoch changes. & Update-path geometry (\cref{fig:path-geometry}), hybrid wrapper equations, recovery claim, and implementation validation of path processing (\cref{tab:evidence}). \\
G2 & Policy bytes are part of the accepted epoch, not advisory application metadata. & Group-context definition (\cref{eq:gch}), epoch key schedule (\cref{eq:epoch-keys}), policy-consensus claim, and \Shield{} flow (\cref{fig:shield}). \\
G3 & \Shield{} selects stricter processing rules through explicit protocol inputs. & Shield-mode construction, relay/member split (\cref{tab:relay}), external-join policy handling, and validation evidence for policy rejection. \\
G4 & Sender chains keep within-epoch traffic efficient without erasing the epoch boundary. & Message algorithm (\cref{alg:message}), sender-chain derivation, replay claim, and message-layer validation. \\
G5 & Franking, sealed content, and expiration are checked at accepted message boundaries. & Message-envelope figure (\cref{fig:message-envelope}), franking figure (\cref{fig:franking-flow}), franking claim, and ProVerif query Q6. \\
G6 & Symbolic and executable evidence support a stated scope rather than a complete group proof. & Claim-boundary table (\cref{tab:scope}), verification-scope table (\cref{tab:group-verification-scope}), ProVerif results (\cref{tab:group-proverif}), and reproducibility notes. \\
\bottomrule
\end{tabularx}
\end{table}

\section{Related Work and Design Position}\label{sec:related}

\subsection{MLS and TreeKEM}

\MLS{} is the closest standard reference point for the group-state problem.
It separates group membership, group-epoch transitions, authenticated group context,
and application-secret derivation into a protocol whose main state transition
is a TreeKEM Commit message~\cite{rfc9420}. Here, a \emph{Commit}
means the MLS/TreeKEM protocol message that applies authenticated Add, Remove,
Update, PSK, or ReInit proposals and installs the next group epoch after
client-side verification. It is a group-state transition message, not an
implementation or repository action. The TreeKEM line of work gives the
asynchronous tree structure that makes large dynamic groups practical: a
member does not need to establish a fresh pairwise secret with every other
member for every update, because refreshed path secrets can be encrypted to
co-path resolutions~\cite{treekem}. TreeSync then sharpens the engineering
problem by treating authenticated tree synchronization as a first-class
security boundary rather than a mere serialization detail~\cite{treesync}.

\Aura{} Group follows this MLS/TreeKEM shape while defining its own protocol
profile rather than an RFC~9420 wire-compatible variant. The design isolates a
technical question central to Aura Messenger's group architecture: how do
hybrid classical/post-quantum update paths and policy-bound group-epoch
derivation affect group-state evolution? This restriction is deliberate. The
paper does not try to replace the MLS standard, and it does not claim the full
security theorem expected for all MLS proposal schedules. It examines a smaller
interface in which the post-quantum contribution is injected at the
membership-changing update path and the application-security policy is bound
before confirmation. That keeps the analyzed cryptographic surface tractable
while preserving the group-specific problems that do not appear in a two-party
ratchet.

\subsection{Sender-Key Messaging and Message Semantics}

Many deployed group messengers use sender-key style techniques for efficient
traffic within a group epoch. A sender key avoids running a full group key update for
every application message, but it shifts attention to per-sender counters,
replay windows, skipped-key storage, sender authentication, and group-epoch
separation. \Aura{} Group adopts this efficiency pattern within a group epoch, but
does not treat it as an independent protocol. Sender chains are derived from
the accepted group-epoch state, and encrypted sender data is tied to the same
metadata key that comes from the policy-bound key schedule.

The message-semantics layer is similarly integrated. Message franking provides
selective accountable disclosure without giving the relay plaintext access to
ordinary traffic~\cite{franking}. Sealed content follows the same design
instinct as sealed-sender metadata minimization: expose only what routing and
delivery need, and keep content-level secrets inside end-to-end encrypted
payloads~\cite{sealed-sender}. In \Aura{} Group these mechanisms are not
treated as product features attached after encryption. They are explicit
application-frame semantics whose tags, timestamps, and references are checked
at the cryptographic processing boundary.

\subsection{Post-Quantum Migration Boundary}

The post-quantum migration pressure is different for pairwise handshakes and
group-epoch intervals. A pairwise protocol can add a KEM encapsulation to the initial
setup or to directed ratchet steps. A group protocol must decide how KEM
material enters a tree update, how the corresponding ciphertexts are addressed
to a co-path resolution, and how the accepted group-epoch interval records that policy and
tree state agreed. \Aura{} Group therefore moves ML-KEM material into the
TreeKEM update path rather than limiting it to account admission or group
creation.

The design remains conservative about authentication. Ed25519 is used for key
packages, Commit messages, and message signatures. That is a classical authentication
boundary and is stated as a limitation rather than hidden behind the
confidentiality claims. The contribution here is hybrid
group-epoch confidentiality and policy agreement for group state; replacing Ed25519
with a post-quantum signature scheme is a separate design and deployment
decision.

\subsection{Two-Party Aura Protocol and Account Admission}

The two-party \Aura{} Protocol manuscript studies a pairwise handshake and
directed ratchet boundary~\cite{melnychenko2026aura}. The Hybrid PQ-OPAQUE
manuscript and its formal-evidence report study password-authenticated account
admission in the augmented-PAKE tradition of Jarecki--Krawczyk--Xu and the
IRTF OPAQUE specification~\cite{melnychenko2026opaque,
melnychenko2026opaqueformal,jarecki2018,ietf-opaque}. \Aura{} Group uses those
layers as surrounding protocol context, not as substitutes for group-security
analysis. After credentials and identities are authenticated, the group
protocol must still define accepted tree state, epoch transitions,
sender-chain traffic, replay control, and relay-visible encrypted envelopes.

\section{Threat Model and Scope}\label{sec:threat}

The adversary controls the network, reorders and replays messages, corrupts
old serialized state, tampers with Commit messages and application ciphertexts, and may
learn secrets of removed members. The relay is not trusted with plaintext or
group secrets. It may validate structure and route messages, but authoritative
membership and accepted group-epoch state are client-side decisions.

The model is single-device. Identity authentication uses Ed25519 signatures
over key packages and Commit messages. This is classical authentication; a
transition to post-quantum signatures is orthogonal to the TreeKEM and
policy-binding mechanisms. The security account combines a scoped \ProVerif{}
model, reference-implementation evidence, and a protocol argument. Complete
theorems for asynchronous concurrent groups, including delivery schedules and
all interleavings of Add, Remove, Update, PSK, and ReInit proposals, are
outside this model.

The adversary may be adaptive at the level of endpoint state snapshots. A
snapshot can reveal stored group-epoch secrets, sender-chain state, replay-window
state, pending proposals, and serialized group state that was not sealed or
has already been restored. The model still assumes correct local cryptographic
erasure after successful group-epoch processing: a member that keeps obsolete path
secrets, group-epoch secrets, or message keys indefinitely has stepped outside the
claim boundary. This is the same practical boundary that appears in pairwise
ratchet protocols, but group messaging makes the state surface larger because
there are per-sender chains and tree nodes in addition to a single pairwise
root key.

The relay is fully adversarial for confidentiality and membership authority.
It may drop a Commit message, delay a Welcome, replay an old message, route one member
through stale public state, or present different clients with different
delivery schedules. Clients therefore accept only local state transitions that
verify the relevant signatures, tree hashes, policy bytes, confirmation MACs,
and replay windows. The relay can still perform useful structural checks,
including size limits and sender binding checks, but those checks are not a
substitute for client-side cryptographic acceptance.

\subsection{Security Goals}

The construction is organized around six goals. First, accepted group-epoch secrets
are intended to be pseudorandom to an adversary that did not learn both the
surviving hybrid path material and the previous init contribution. Second, an
honest post-compromise update creates a recovery boundary at the next accepted
group epoch, provided at least one hybrid component remains unbroken and old secrets
are erased. Third, removed or not-yet-added members are excluded from messages
outside their membership interval. Fourth, members that disagree on policy fail
to share a group epoch silently. Fifth, application messages reject replay inside
the accepted group epoch. Sixth, frankable messages expose only selective
disclosures to a moderator.

\begin{table}[tbp]
\centering
\caption{Claim boundary for the group paper.}\label{tab:scope}
\footnotesize
\begin{tabularx}{\textwidth}{@{}p{0.29\textwidth}YY@{}}
\toprule
Topic & Claimed here & Not claimed here \\
\midrule
Group-epoch secrecy & Fresh group-epoch keys after accepted hybrid TreeKEM Commit messages, subject to at least one unbroken component and honest erasure. & Universal-composability proof for all concurrent group schedules. \\
Post-compromise recovery & Recovery at group-epoch boundaries when an honest update path introduces fresh \DH{} and \KEM{} material after compromise. & Recovery without a later honest group-epoch update. \\
Policy enforcement & Policy bytes are hashed into the group context and affect group-epoch keys and confirmation MACs. & Server-side policy authority or mutable policy negotiation. \\
Relay role & Relay-blind structural validation and routing. & Relay knowledge of plaintext, group secrets, or authoritative membership state. \\
Application semantics & Franking, sealed content, and expiration are enforced at message processing boundaries. & Legal moderation semantics beyond cryptographic verification of a disclosed message. \\
\bottomrule
\end{tabularx}
\end{table}

\section{Construction Overview}

\Aura{} Group has three layers. The group-epoch layer maintains a ratchet tree and
advances it with Commit messages. The key-schedule layer derives group-epoch secrets,
metadata keys, welcome keys, confirmation keys, and sender-key bases. The
message layer encrypts application payloads under per-sender chains while
encrypting sender metadata under the group-epoch metadata key.

The tree is left-balanced and stores both \DH{} and \KEM{} public material at
leaves and internal nodes~\cite{rfc7748,fips203}. A member's key package binds
identity \EdDSA{},
identity \DH{}, identity \KEM{}, leaf \DH{}, leaf \KEM{}, and credential bytes
under an \EdDSA{} signature. This separates long-term identity from group-epoch
update material while giving recipients a signed object for membership
admission.

\begin{figure}[H]
\centering
\resizebox{0.985\textwidth}{!}{%
\begin{tikzpicture}[
  font=\small,
  node distance=0.75cm,
  block/.style={draw,rounded corners=2pt,align=center,minimum width=2.75cm,
    minimum height=0.86cm,fill=blue!5},
  smallblock/.style={draw,rounded corners=2pt,align=center,minimum width=2.45cm,
    minimum height=0.78cm,fill=green!6},
  arrow/.style={-Latex,thick}
]
\node[block] (tree) at (0,0) {Ratchet tree\\leaf and path keys};
\node[block] (commit) at (3.55,0) {Hybrid TreeKEM\\UpdatePath};
\node[block] (ctx) at (7.1,0) {Policy-bound\\group context};
\node[block] (epoch) at (11.05,0) {Group-epoch keys\\metadata, welcome, confirm};
\node[smallblock] (sender) at (11.05,-1.55) {Sender chains\\per leaf};
\node[smallblock] (msg) at (14.4,-1.55) {Application\\messages};

\draw[arrow] (tree) -- (commit);
\draw[arrow] (commit) -- (ctx);
\draw[arrow] (ctx) -- (epoch);
\draw[arrow] (epoch.south) -- (sender.north);
\draw[arrow] (sender) -- (msg);
\end{tikzpicture}%
}
\caption{Group-epoch transition and message-key flow. A Commit message updates the tree,
derives a policy-bound context, and refreshes sender chains for the new
membership interval.}
\label{fig:epoch-flow}
\end{figure}
\FloatBarrier

\subsection{Key Packages and Welcome}

Adding a member produces a signed key package, an Add Commit message, and a Welcome message.
The Welcome encrypts the joiner secret for the added member and carries enough
public group information for the recipient to reconstruct the tree, verify the
tree hash, parse the serialized policy, derive group-epoch keys, and check the
confirmation MAC. The design therefore keeps admission state explicit: a
member joins after decrypting Welcome and accepting the resulting group-epoch state,
not because the relay appended an identity to a roster.

\begin{figure}[H]
\centering
\resizebox{\textwidth}{!}{%
\begin{tikzpicture}[
  font=\small,
  joinstep/.style={draw,rounded corners=2pt,align=center,minimum width=3.0cm,minimum height=0.78cm,fill=blue!6},
  check/.style={draw,rounded corners=2pt,align=center,minimum width=2.75cm,minimum height=0.68cm,fill=green!7,font=\scriptsize},
  arrow/.style={-Latex,thick},
  node distance=0.62cm and 0.74cm
]
\node[joinstep] (kp) {Joiner\\signed KeyPackage};
\node[joinstep,right=of kp] (commit) {Member creates\\Add Commit};
\node[joinstep,right=of commit] (welcome) {Welcome encrypts\\joiner secret};
\node[joinstep,right=of welcome] (rebuild) {Joiner rebuilds\\tree and group epoch};
\node[joinstep,right=of rebuild] (accept) {Accepted\\group state};

\node[check,below=of rebuild] (verify) {verify tree hash\\policy bytes\\confirmation MAC};
\node[check,below=of welcome] (secret) {decrypt Welcome\\with leaf DH+KEM};

\draw[arrow] (kp) -- (commit);
\draw[arrow] (commit) -- (welcome);
\draw[arrow] (welcome) -- (rebuild);
\draw[arrow] (rebuild) -- (accept);
\draw[arrow] (secret) -- (welcome);
\draw[arrow] (verify) -- (rebuild);
\end{tikzpicture}%
}
\caption{Welcome/admission boundary. A new member becomes authoritative only after reconstructing the tree, parsing the policy, deriving the group-epoch keys, and verifying the confirmation MAC.}\label{fig:welcome-admission}
\end{figure}

\begin{table}[H]
\centering
\caption{Group-state objects and their role.}\label{tab:objects}
\footnotesize
\begin{tabularx}{\textwidth}{@{}p{0.24\textwidth}YY@{}}
\toprule
Object & Contents & Role in the construction \\
\midrule
Key package & Identity keys, leaf update keys, credential, signature. & Authenticated input for adding a member to the tree. \\
UpdatePath & Leaf public keys, path node public keys, encrypted path secrets, parent hashes. & Refreshes group-epoch tree material for current members. \\
Group context hash & Group identifier, group-epoch number, tree hash, serialized policy. & Binds key derivation to tree state and security policy. \\
Welcome & Encrypted joiner secret and public group information. & Allows a newly added member to reconstruct and verify group-epoch state. \\
Application message & Encrypted sender data, encrypted payload, nonces, signature, optional franking tag. & Carries traffic within a group epoch under sender-chain keys. \\
\bottomrule
\end{tabularx}
\end{table}

\subsection{Tree Path Geometry}

The construction follows the TreeKEM intuition that a member refreshes the
nodes on its direct path and encrypts the corresponding path secrets to the
co-path resolution. \Cref{fig:path-geometry} shows the case of leaf $2$ in an
eight-member tree. The shaded path is recomputed by the committer; the marked
co-path nodes are sufficient recipients for the encrypted path secrets. The
figure also shows the operational effect of the Commit: fresh path entropy is
expanded along the direct path, wrapped to the co-path resolution, and then
used as the commit secret for the next group epoch.

\begin{figure}[!tbp]
\centering
\makebox[\textwidth][c]{%
\resizebox{1.08\textwidth}{!}{%
\begin{tikzpicture}[
  font=\small,
  leaf/.style={draw,circle,minimum size=1.00cm,inner sep=1.2pt,fill=gray!8},
  actor/.style={draw,circle,minimum size=1.18cm,inner sep=1.2pt,fill=blue!18,thick},
  inode/.style={draw,circle,minimum size=0.94cm,inner sep=1pt,fill=gray!5},
  pathnode/.style={draw,circle,minimum size=1.18cm,inner sep=1.2pt,fill=blue!16,thick,font=\footnotesize},
  copath/.style={draw,circle,minimum size=1.18cm,inner sep=1.2pt,fill=orange!18,thick,font=\footnotesize},
  edge/.style={thick,gray!58},
  pathedge/.style={line width=1.7pt,blue!68!black},
  wrap/.style={-Latex,thick,orange!85!black},
  flow/.style={draw,rounded corners=2pt,align=center,
    inner sep=4.5pt,font=\scriptsize,text width=2.28cm,minimum height=0.92cm},
  flowarrow/.style={-Latex,thick,blue!60!black},
  weakflow/.style={-Latex,thick,dashed,blue!60!black},
  legend/.style={font=\scriptsize,align=left}
]

\foreach \i in {0,1,2,3,4,5,6,7} {
  \coordinate (l\i) at (1.55*\i,0);
}
\coordinate (n01) at (0.775,1.18);
\coordinate (n23) at (3.875,1.18);
\coordinate (n45) at (6.975,1.18);
\coordinate (n67) at (10.075,1.18);
\coordinate (n03) at (2.325,2.58);
\coordinate (n47) at (8.525,2.58);
\coordinate (n07) at (5.425,4.10);

\draw[edge] (l0) -- (n01) -- (l1);
\draw[edge] (l2) -- (n23) -- (l3);
\draw[edge] (l4) -- (n45) -- (l5);
\draw[edge] (l6) -- (n67) -- (l7);
\draw[edge] (n01) -- (n03) -- (n23);
\draw[edge] (n45) -- (n47) -- (n67);
\draw[edge] (n03) -- (n07) -- (n47);

\draw[pathedge] (l2) -- (n23) -- (n03) -- (n07);

\foreach \i in {0,1,4,5,6,7} {
  \node[leaf] at (l\i) {$\i$};
}
\node[actor] (committerLeaf) at (l2) {$2$};
\node[copath] (cp3) at (l3) {$3$};
\node[copath] (cpr1) at (n01) {$r_1$};
\node[inode] at (n45) {};
\node[inode] at (n67) {};
\node[copath] (cpr2) at (n47) {$r_2$};
\node[pathnode] (path0) at (n23) {$p_0$};
\node[pathnode] (path1) at (n03) {$p_1$};
\node[pathnode] (path2) at (n07) {$p_2$};

\node[font=\scriptsize,blue!70!black,align=center,fill=white,inner sep=1.2pt] at (2.98,-1.18)
  {committer\\leaf};
\node[font=\scriptsize,blue!70!black,align=center,fill=white,inner sep=1.2pt] at (5.42,5.32)
  {refreshed\\direct path};

\draw[wrap,bend left=24] (path0.east) to node[pos=0.35,right,font=\scriptsize,fill=white,inner sep=1.2pt]
  {$C_{\mathsf{path},0}$} (cp3.north west);
\draw[wrap,bend right=22] (path1.west) to node[pos=0.38,left,font=\scriptsize,fill=white,inner sep=1.2pt]
  {$C_{\mathsf{path},1}$} (cpr1.north east);
\draw[wrap,bend left=22] (path2.east) to node[pos=0.38,above,font=\scriptsize,fill=white,inner sep=1.2pt]
  {$C_{\mathsf{path},2}$} (cpr2.north west);

\node[legend,anchor=north west] at (8.70,4.87)
  {blue path: refreshed secrets\\orange nodes: co-path recipients\\orange arrows: hybrid-wrapped path secrets};

\node[flow,fill=blue!6] (seed) at (1.48,-2.42)
  {\textbf{1. Fresh entropy}\\sample path seed};
\node[flow,fill=blue!6,right=0.42cm of seed] (derive)
  {\textbf{2. Path derivation}\\derive $p_0,p_1,p_2$};
\node[flow,fill=orange!8,right=0.42cm of derive] (wraps)
  {\textbf{3. Hybrid wrapping}\\to $\{3,r_1,r_2\}$};
\node[flow,fill=green!7,right=0.42cm of wraps] (keys)
  {\textbf{4. Epoch keys}\\$\commitsec_{e+1}$ drives\\$e+1$ keys};

\draw[weakflow] (committerLeaf.south west) -- node[pos=0.40,left,font=\scriptsize,fill=white,inner sep=1.2pt]
  {Commit operation} (seed.north);
\draw[flowarrow] (seed) -- (derive);
\draw[flowarrow] (derive) -- (wraps);
\draw[flowarrow] (wraps) -- (keys);
\end{tikzpicture}%
}%
}
\caption{Core Hybrid TreeKEM UpdatePath and group-epoch effect for a Commit from
leaf $2$. The committer refreshes its direct path, encrypts the new path
secrets to the co-path resolution, and feeds the recovered commit secret into
the next group-epoch key schedule.}
\label{fig:path-geometry}
\end{figure}

\subsection{Hybrid Update Path}

For a Commit at group epoch $e$, the committer samples a fresh path secret and
derives node secrets along its direct path. Each updated node receives a fresh
\DH{} keypair and a fresh \KEM{} keypair. In this profile, Add, Remove, and
Update transitions that refresh tree secrets carry an UpdatePath; PSK and
ReInit transitions still advance the epoch state but are not counted in the
one-step hybrid path-secrecy argument unless they also carry fresh path
material. A co-path resolution is the TreeKEM recipient set covering sibling
subtrees that need the refreshed path secret in order to reconstruct the new
tree state. For every recipient resolution node $R$, the path secret is wrapped
under a hybrid key:
\begingroup
\footnotesize
\begin{equation*}
\begin{aligned}
  (x_{\mathsf{eph}},X_{\mathsf{eph}})&\gets \DH.\KeyGen(),\\
  z_{\DH}&=\DH(x_{\mathsf{eph}},X_R),\\
  (c_{\KEM},z_{\KEM})&\gets \KEM.\Encaps(pk_R),\\
  s_{\mathsf{path}}&=\domain{Aura-PQ-Group-Hybrid}\concat z_{\DH},\\
  i_{\mathsf{path}}&=\domain{Aura-Group-HybridPath},\\
  k_{\mathsf{path}}&=\HKDF(z_{\KEM},32,s_{\mathsf{path}},i_{\mathsf{path}}),\\
  \mathsf{aad}_{\mathsf{path}}&=\gid\concat e\concat \mathit{nodeIndex}
    \concat X_{\mathsf{eph}}\concat c_{\KEM},\\
  C_{\mathsf{path}}&=\AEAD.\Enc(k_{\mathsf{path}},n_{\mathsf{path}},
    \mathit{pathSecret},\mathsf{aad}_{\mathsf{path}}).
\end{aligned}
\end{equation*}
\endgroup
Here $\HKDF(x,L,s,i)$ denotes HKDF-SHA256 with input keying material $x$,
output length $L$, salt $s$, and domain-separation information $i$. This
matches the construction's hybrid path wrapper: the \KEM{} shared secret is
the input material, and the \DH{} output is carried in the salt with the
path-domain label.
The corresponding recipient decapsulates $c_{\KEM}$, computes the matching
\DH{} value, reconstructs $k_{\mathsf{path}}$, and decrypts the path secret.
The group identifier, group-epoch number, node index, ephemeral \DH{} public key, and \KEM{}
ciphertext are associated data, so a wrapped path secret cannot be moved to a
different group, group-epoch interval, tree position, or hybrid wrapper without detection.
Parent hashes are verified along the updated path so that an adversary cannot
splice a valid node into an inconsistent tree transcript. HKDF and the AEAD
interface follow the standard extract-then-expand and nonce-misuse-resistant
AEAD references~\cite{rfc5869,rfc8452}.

\subsection{Policy-Bound Group-Epoch Context}

A \emph{group epoch} in \Aura{} Group is a cryptographic membership interval:
one accepted ratchet tree, one serialized policy, one group-epoch secret, one
metadata key, and one set of per-member sender chains. Group epoch $0$ is the
creator-only genesis interval. Every accepted Add, Remove, Update, PSK, or
ReInit Commit message advances the group from $e$ to $e+1$ atomically; application
messages only consume sender-chain generations inside the current interval,
and a rejected Commit leaves local state unchanged. The transition occurs only
at an authenticated state boundary: the accepted TreeKEM Commit message, where
tree hash, policy bytes, confirmation MAC, and group-epoch number agree before
new keys are installed.

The terminology is fixed as follows. The \emph{epoch number} $e$ is the
monotonic counter carried in group messages and Commit messages. The
\emph{group context hash} $\gch_e$ is the hash of the group identifier, epoch
number, tree hash, and serialized policy. The \emph{epoch secret}
$\epochsec_e$ is derived only after an accepted transition into epoch $e$. The
\emph{accepted epoch state} is the local installed tuple consisting of the
tree, policy, context hash, epoch keys, sender-chain bases, replay state, and
sealed-state freshness metadata. Below, ``group epoch'' denotes the membership
interval, and ``accepted epoch state'' denotes the installed local state.

The group-epoch context is defined with length prefixes for variable-size fields:
\begin{equation}
\begin{aligned}
  \gch_e =
  \SHA\!\bigl(&
    \ell_{32}(\gid)\concat \gid \concat
    \mathrm{le}_{64}(e)\\
    &\concat\ell_{32}(\treehash_e)\concat \treehash_e
    \concat\ell_{32}(\policy)\concat \policy
  \bigr).
\end{aligned}
  \label{eq:gch}
\end{equation}
Here $\ell_{32}(x)$ is the 32-bit little-endian length of $x$, and
$\mathrm{le}_{64}(e)$ is the 64-bit little-endian group-epoch encoding. The length
prefixes remove concatenation ambiguity. More importantly, the serialized
policy is inside the hash. If two members disagree on whether
\Shield{} is enabled, or on the sender-key limits, they compute different
$\gch_e$, different group-epoch keys, and different confirmation MACs.
For compact pseudocode, write
$\mathsf{GroupContext}(\gid,e,\treehash,\policy)$ for the hash in
\cref{eq:gch}.

For an UpdatePath-bearing accepted transition into group epoch $e$, let
$\commitsec_e$ be the fresh secret recovered from the accepted UpdatePath and
let $\initsec_{e-1}$ be the chaining secret carried from the previous group
epoch. PSK-only or ReInit-only transitions may still advance the epoch state
through their own transition secret, but they are outside the one-step hybrid
path-secrecy argument unless they also carry fresh path material. The key
schedule first combines continuity with transition entropy and then binds the
result to the policy-bound group context:
\[
  \mathcal{A}=\{\mathsf{meta},\mathsf{welcome},\mathsf{confirm},\mathsf{init}\}.
\]
\begin{align}
  j_e &= \HKDF\text{-}\mathrm{Extract}(\initsec_{e-1}, \commitsec_e),\\
  \epochsec_e &=
  \HKDF\text{-}\mathrm{Expand}(j_e,32,\;
    \domain{Aura-Group-EpochSecret}\concat \gch_e),\\
  k^\alpha_e &=
  \mathsf{KDF}^{\policy}_\alpha(\epochsec_e),
  \qquad \alpha\in\mathcal{A},\\
  \mathbf{k}_e &:=
  (k^{\mathsf{meta}}_e,k^{\mathsf{welcome}}_e,
   k^{\mathsf{confirm}}_e,\initsec_e).
  \label{eq:epoch-keys}
\end{align}
Here $\mathsf{KDF}^{\policy}_\alpha$ is selected by the serialized policy.
For the default profile,
\[
  \mathsf{KDF}^{\policy}_\alpha(x)=
  \HKDF\text{-}\mathrm{Expand}(x,32,\mathsf{label}(\alpha)).
\]
The default label map is:
\[
\begin{array}{rcl}
  \mathsf{label}(\mathsf{meta}) &=& \domain{Aura-Group-MetadataKey},\\
  \mathsf{label}(\mathsf{welcome}) &=& \domain{Aura-Group-WelcomeKey},\\
  \mathsf{label}(\mathsf{confirm}) &=& \domain{Aura-Group-ConfirmKey},\\
  \mathsf{label}(\mathsf{init}) &=& \domain{Aura-Group-InitSecret}.
\end{array}
\]
Write $\mathrm{Subkeys}_{\policy}(x)$ for the tuple
$(k^{\mathsf{meta}},k^{\mathsf{welcome}},k^{\mathsf{confirm}},\initsec)$
derived by these policy-selected KDFs.

The four outputs have separate protocol roles.
$k^{\mathsf{meta}}_e$ encrypts sender data,
$k^{\mathsf{welcome}}_e$ domain-separates Welcome-admission material and
supplemental Welcome fields,
$k^{\mathsf{confirm}}_e$ authenticates the accepted group context, and
$\initsec_e$ is the chaining secret for the next group-epoch transition.
The extract output $j_e$ is the group-epoch joiner secret; in Add transitions
the core joiner secret is sealed to the joining member's key package inside a
Welcome rather than encrypted under a group-wide key. A recipient recomputes the
group-epoch secret $\epochsec_e$ only after reconstructing the same
$(\gid,e,\treehash_e,\policy)$ context.
The corresponding confirmation tag is:
\begin{equation}
  \tau_e = \HMAC(k^{\mathsf{confirm}}_e,\gch_e).
  \label{eq:confirmation}
\end{equation}
Thus the confirmation step checks tree state and policy state together.

\subsection{Group-Epoch Installation and State Continuity}

Group-epoch installation is transactional. A client first builds a candidate tree,
candidate group context, candidate group-epoch secret, candidate subkeys, and
candidate sender-chain bases. None of these candidate values become live state
until the Commit signature, proposal validity, parent hashes, context binding,
and confirmation MAC have all verified. For malformed input or a failed
cryptographic check, the state transition is pure rejection:
\begin{equation}
  \mathsf{ProcessCommit}(S,\mathsf{com})=(\bot,S).
  \label{eq:group-reject-no-state}
\end{equation}
For accepted Commit messages, the abstract transition is a single step
$S_e\rightarrow S_{e+1}$, even if an implementation parses, decrypts, derives,
persists, notifies callbacks, and erases temporary secrets in separate
instructions. Plaintext or membership callbacks should be released only after
the corresponding replay state, group-epoch counter, sender-chain state, and sealed
state have reached an equivalent durable boundary:
\begin{equation}
  \mathsf{Accept}(S_e,\mathsf{com})=S_{e+1}
  \quad\Longrightarrow\quad
  \mathsf{Persist}(S_{e+1})\prec \mathsf{Expose}(\mathrm{effect}).
  \label{eq:group-atomic-accept}
\end{equation}
This is not an additional cryptographic assumption; it is the state-continuity
condition needed for the replay and membership claims to hold. A crash between
plaintext release and replay-window persistence would
violate the protocol state premise even if all AEAD tags and signatures were
valid.

\begin{table}[tbp]
\centering
\caption{Group-epoch state transitions and non-transitions.}\label{tab:epoch-lifecycle}
\small
\begin{tabularx}{\textwidth}{@{}p{0.22\textwidth}YY@{}}
\toprule
Event & Effect on group-epoch state & Security consequence \\
\midrule
Genesis creation & Creates group epoch $0$ from the initial tree, policy, and group context. & Establishes the first authoritative policy-bound state. \\
Add, Remove, Update, PSK, ReInit Commit message & Advances $e$ to $e+1$ only after signature, proposal, applicable path, parent-hash, context, and confirmation checks. & Refreshes tree secrets when path material is present, resets sender-chain bases, and defines the next membership interval. \\
Application message & Does not advance the group-epoch counter. It consumes a sender-chain generation inside the current membership interval. & Preserves low-cost traffic while keeping replay state sender-indexed and group-epoch-bound. \\
Delayed message from a previous group epoch & May be accepted only from bounded compatibility state tagged by group epoch. & Does not reinstall an old tree or revive old Commit authority. \\
Rejected Commit message or ciphertext & Leaves the live group epoch, tree, sender chains, replay windows, and sealed-state freshness unchanged. & Prevents malformed input from creating partial state transitions. \\
Re-addition after removal & Requires a later Add or valid external join that creates a new membership interval. & Does not grant access to messages from the absence interval. \\
\bottomrule
\end{tabularx}
\end{table}

\subsection{Shield Mode}

\Shield{} is a fixed stricter policy profile. Its default values are profile
parameters rather than theorem parameters: a $1000$-generation limit per sender
chain within a group epoch, $4$ skipped sender keys per sender, external joins
blocked, profile-specific key schedule enabled, and mandatory franking enabled.
The bounds are chosen to limit sender-chain lifetime and skipped-key storage in
high-risk groups; the security statements depend on enforcement of the bounds,
not on the exact numerical constants. The \Shield{} subkey schedule uses two
HKDF-Expand passes with distinct labels:
\begin{align}
  t_\alpha &=
    \HKDF\text{-}\mathrm{Expand}(\epochsec_e,32,
      \mathsf{label}(\alpha)\concat\domain{Aura-Enhanced-Pass1}),\\
  k^\alpha_e &=
    \HKDF\text{-}\mathrm{Expand}(t_\alpha,32,
      \mathsf{label}(\alpha)\concat\domain{Aura-Enhanced-Pass2}).
\end{align}
The second pass is not a claim of extra strength beyond HKDF. Its role is
profile selection and domain separation: members that disagree on the
\Shield{} policy compute a different key schedule and fail confirmation
instead of silently sharing a downgraded epoch.

\begin{figure}[tbp]
\centering
\begin{tikzpicture}[
  font=\small,
  box/.style={draw,rounded corners=2pt,align=center,minimum width=3.2cm,minimum height=0.78cm,fill=orange!8},
  arrow/.style={-Latex,thick},
  node distance=0.9cm and 1.2cm
]
\node[box] (policy) {Serialized\\security policy};
\node[box,right=of policy] (hash) {Group context\\hash};
\node[box,right=of hash] (keys) {Group-epoch keys};
\node[box,below=of keys] (confirm) {Confirmation\\MAC};
\node[box,below=of hash] (sender) {Sender-key\\limits};
\node[box,below=of policy] (frank) {Mandatory\\franking};
\node[box,below=of sender] (join) {External join\\blocked};

\draw[arrow] (policy) -- (hash);
\draw[arrow] (hash) -- (keys);
\draw[arrow] (keys) -- (confirm);
\draw[arrow] (policy) -- (frank);
\draw[arrow] (policy) -- (sender);
\draw[arrow] (policy) -- (join);
\draw[arrow] (sender) -- (keys);
\end{tikzpicture}
\caption{\Shield{} as a cryptographic policy profile. Policy bytes affect
group-epoch derivation and also drive local message-processing limits.}
\label{fig:shield}
\end{figure}

\section{Protocol Operations}\label{sec:operations}

This section gives state-transition pseudocode for the construction. It names
the inputs that must be authenticated, the state fields that must be updated,
and the checks that cause rejection. It is not a wire-format grammar; the
operations are specified at the cryptographic acceptance boundary.

\subsection{Group Creation and Admission}

A group begins with a creator leaf, an initial policy, and a ratchet tree that
contains the creator's signed leaf material. Adding a member is a group-epoch
transition, not a roster mutation performed by the relay. The new member must
receive a Welcome, decrypt the joiner secret, rebuild the tree, verify the
tree hash, parse the policy, derive group-epoch keys, and check the confirmation MAC
before it treats itself as a member.

\begingroup
\LinesNotNumbered
\SetAlgoNoLine
\begin{algorithm}[H]
\caption{Create a policy-bound group and add one member.}\label{alg:create-add}
\footnotesize
\textbf{Inputs.}
Creator identity $I_A$, joiner key package $KP_B$, and serialized policy $\policy$.

\textbf{Outputs.}
Updated creator state $S_A'$, Add Commit message $\mathsf{com}$, and Welcome message $\mathsf{wel}$.
\vspace{0.25\baselineskip}

\textbf{1. Create the genesis group state.}
The creator builds the initial ratchet tree and derives the genesis context:
\[
\begin{aligned}
S_A &\gets \mathrm{CreateTree}(I_A,\policy),\\
\gid &\gets S_A.\gid,\qquad
\treehash_0\gets \mathsf{TreeHash}(S_A.\mathrm{tree}),\\
\gch_0 &\gets \mathsf{GroupContext}(\gid,0,\treehash_0,\policy),\\
(\epochsec_0,\initsec_0)&\gets \mathrm{InitEpoch}(\gch_0).
\end{aligned}
\]

\textbf{2. Validate and insert the joining leaf.}
Reject unless $KP_B$ verifies as a signed key package. If it verifies, insert it
into the tree and compute the hybrid update path:
\[
\begin{aligned}
T_1 &\gets \mathrm{InsertLeaf}(S_A.\mathrm{tree},KP_B),\\
(T_1,\mathsf{path},\commitsec_1)&\gets
  \mathrm{HybridUpdatePath}(T_1,S_A.\mathrm{leaf}).
\end{aligned}
\]

\textbf{3. Derive the next group-epoch state.}
\[
\begin{aligned}
\treehash_1&\gets \mathsf{TreeHash}(T_1),\\
\gch_1&\gets \mathsf{GroupContext}(\gid,1,\treehash_1,\policy),\\
j_1&\gets
  \HKDF\text{-}\mathrm{Extract}(\initsec_0,\commitsec_1),\\
\epochsec_1&\gets
  \HKDF\text{-}\mathrm{Expand}(j_1,32,
  \domain{Aura-Group-EpochSecret}\concat\gch_1),\\
\left(\begin{gathered}
  k^{\mathsf{meta}}_1,k^{\mathsf{welcome}}_1\\
  k^{\mathsf{confirm}}_1,\initsec_1
\end{gathered}\right)
  &\gets \mathrm{Subkeys}_{\policy}(\epochsec_1),\\
\tau_1&\gets \HMAC(k^{\mathsf{confirm}}_1,\gch_1).
\end{aligned}
\]

\textbf{4. Produce the authenticated admission messages.}
\[
\begin{aligned}
\mathsf{com}&\gets
  \mathrm{Sign}_{I_A}(\gid,0,1,\mathsf{Add}(KP_B),\mathsf{path},\gch_1,\tau_1),\\
\mathsf{wel}&\gets
  \mathrm{SealWelcome}_{KP_B}
    (\gid,1,j_1,T_1,\treehash_1,\policy,\tau_1).
\end{aligned}
\]

\textbf{5. Install the creator's accepted group epoch.}
\[
\begin{aligned}
S_A'&\gets \mathrm{InstallEpoch}(S_A,T_1,1,\policy,\epochsec_1,\initsec_1),\\
S_A'.\mathrm{senderChains}&\gets
  \mathrm{ResetSenderChains}(\epochsec_1,S_A'.\mathrm{members}).
\end{aligned}
\]
Return $(S_A',\mathsf{com},\mathsf{wel})$.
\end{algorithm}
\endgroup

\begingroup
\LinesNotNumbered
\SetAlgoNoLine
\begin{algorithm}[H]
\caption{Process a Welcome as the joining member.}\label{alg:welcome}
\footnotesize
\textbf{Inputs.}
Welcome message $\mathsf{wel}$, local joiner secret-key material, and expected identity $I_B$.

\textbf{Output.}
Accepted member state $S_B$, or rejection.
\vspace{0.25\baselineskip}

\textbf{1. Open Welcome and parse policy.}
The joiner decrypts $\mathsf{wel}$ under its leaf credentials. It rejects if
opening fails or if the serialized policy cannot be parsed.

\textbf{2. Reconstruct the public group state.}
From the Welcome information, the joiner rebuilds the ratchet tree $T$ and
checks that $I_B$ has a leaf, that the roster verifies, and that
\[
  \mathsf{TreeHash}(T)=\mathsf{info.treeHash}.
\]

\textbf{3. Recompute the policy-bound context and subkeys.}
\[
\begin{aligned}
\gch&\gets \mathsf{GroupContext}(\mathsf{info.groupId},
  \mathsf{info.epoch},\mathsf{TreeHash}(T),\policy),\\
\epochsec&\gets
  \mathrm{DeriveEpochSecret}(\mathsf{info.joinerSecret},\gch),\\
\left(\begin{gathered}
  k^{\mathsf{meta}},k^{\mathsf{welcome}}\\
  k^{\mathsf{confirm}},\initsec
\end{gathered}\right)
  &\gets \mathrm{Subkeys}_{\policy}(\epochsec).
\end{aligned}
\]

\textbf{4. Verify confirmation and install state.}
Reject unless
\[
  \HMAC(k^{\mathsf{confirm}},\gch)=\mathsf{info.confirmation}.
\]
On success, install the accepted group epoch and reset sender chains:
\[
S_B\gets \mathrm{InstallEpoch}(T,\mathsf{info.groupId},
  \mathsf{info.epoch},\policy,\epochsec,\initsec).
\]
The sender-chain bases of $S_B$ are reset from $\epochsec$ for the accepted roster.
Return $S_B$.
\end{algorithm}
\endgroup
\FloatBarrier

\subsection{Group-Epoch Transition Processing}

An Update, Add, Remove, PSK, or ReInit proposal becomes authoritative only
after a signed Commit message is accepted. For the UpdatePath-bearing Commit
messages analyzed here, the Commit identifies the old group epoch, the new group
epoch, the proposal set, the committer leaf, the update path, and the
confirmation MAC. Members process the update path from their own position in the
tree: they locate the first path secret encrypted to a node in their resolution,
decrypt it with the hybrid wrapper, derive the remaining path nodes, verify
parent hashes, and then derive the group-epoch keys. PSK-only and ReInit-only
Commit messages use the same context and confirmation checks but obtain their
transition secret from proposal processing rather than from
$\mathrm{DecryptOrApplyUpdatePath}$.

\begingroup
\LinesNotNumbered
\SetAlgoNoLine
\begin{algorithm}[!tbp]
\caption{Process an UpdatePath-bearing group-epoch Commit message.}\label{alg:process-commit}
\footnotesize
\textbf{Inputs.}
Local state $S$ and signed Commit message $\mathsf{com}$.

\textbf{Output.}
New accepted state $S'$, removal notification, or rejection.
\vspace{0.25\baselineskip}

\textbf{1. Check the public transition envelope.}
Reject unless the Commit is for $S.\gid$, advances from $S.\epoch$ to
$S.\epoch+1$, carries the same policy bytes as $S.\policy$, has a valid
committer signature, and contains proposals valid for the current roster.

\textbf{2. Apply public tree changes.}
Apply the public proposal effects to obtain a candidate tree $T$. If the local
leaf is removed, erase local group-epoch secrets and return $\mathrm{Removed}$.

\textbf{3. Recover and authenticate the update path.}
From the local leaf position, decrypt or apply the first reachable path secret:
\[
  (T,\commitsec)\gets
  \mathrm{DecryptOrApplyUpdatePath}(T,\mathsf{com.updatePath},S.\mathrm{leaf}).
\]
Reject if no path secret can be recovered or if the parent hashes do not verify
against $\mathsf{com.committer}$.

\textbf{4. Recompute the accepted group context.}
\[
  \gch\gets \mathsf{GroupContext}(S.\gid,S.\epoch+1,
  \mathsf{TreeHash}(T),S.\policy).
\]
Reject unless $\mathsf{com.context}=\gch$.

\textbf{5. Derive the new group-epoch keys.}
\[
\begin{aligned}
j&\gets \HKDF\text{-}\mathrm{Extract}(S.\initsec,\commitsec),\\
\epochsec&\gets \HKDF\text{-}\mathrm{Expand}(j,32,
  \domain{Aura-Group-EpochSecret}\concat\gch),\\
\left(\begin{gathered}
  k^{\mathsf{meta}},k^{\mathsf{welcome}}\\
  k^{\mathsf{confirm}},\initsec
\end{gathered}\right)
  &\gets \mathrm{Subkeys}_{S.\policy}(\epochsec).
\end{aligned}
\]
Reject unless $\HMAC(k^{\mathsf{confirm}},\gch)=\mathsf{com.confirmation}$.

\textbf{6. Install the new group epoch atomically.}
\[
\begin{aligned}
S'&\gets \mathrm{InstallEpoch}(S,T,S.\epoch+1,S.\policy,\epochsec,\initsec),\\
S'.\mathrm{senderChains}&\gets
  \mathrm{ResetSenderChains}(\epochsec,S'.\mathrm{members}).
\end{aligned}
\]
Return $S'$.
\end{algorithm}
\endgroup
\FloatBarrier

\subsection{Application Messages}

Application messages have lower computational and bandwidth cost than
group-epoch transition Commit messages. The sender advances its
own chain, encrypts sender data under the group-epoch metadata key, encrypts payload
content under the message key, signs the envelope, and optionally attaches a
franking tag. The receiver authenticates the signature, decrypts sender data,
checks the replay window, derives or retrieves the sender-chain key, decrypts
the payload, and then applies policy checks such as mandatory franking,
expiration, maximum size, and allowed message semantics.
The sender and receiver state updates are monotone: a sent message persists the
advanced sender chain, and an accepted receive persists replay-window state
before plaintext is exposed to the application.

\begingroup
\LinesNotNumbered
\SetAlgoNoLine
\begin{algorithm}[!tbp]
\caption{Encrypt and decrypt a group-epoch application message.}\label{alg:message}
\footnotesize
\textbf{Inputs.}
Sender state $S_i$, receiver state $S_j$, and plaintext $m$.

\textbf{Outputs.}
Updated sender state $S_i'$, ciphertext $C$, and either updated receiver state
$S_j'$ with plaintext $m'$ or rejection.
\vspace{0.25\baselineskip}

\textbf{1. Sender advances its chain and prepares sender data.}
\[
\begin{aligned}
(mk,ck_i')&\gets \mathrm{AdvanceSenderChain}(S_i.ck_i),\\
(n_{\mathsf{sd}},n_{\mathsf{pt}},\mathit{reuseGuard})
  &\gets \mathrm{FreshMessageRandomness}(),\\
sd&\gets (i,S_i.\mathrm{generation},\mathit{reuseGuard}),\\
\mathsf{aad}_{\mathsf{sd}}&\gets S_i.\gid\concat S_i.\epoch,\\
C_{\mathsf{sd}}&\gets
  \AEAD.\Enc(S_i.k^{\mathsf{meta}},n_{\mathsf{sd}},sd,\mathsf{aad}_{\mathsf{sd}}).
\end{aligned}
\]

\textbf{2. Sender applies message semantics and encrypts the payload.}
If franking is required or requested, compute $(fk,\mathit{tag})\gets
\mathrm{Frank}(m)$ and set $m^\star\gets(m,fk)$; otherwise set
$\mathit{tag}\gets\varepsilon$ and $m^\star\gets m$. Then compute
\[
\begin{aligned}
\mathsf{aad}_{\mathsf{pt}}&\gets
  \mathsf{aad}_{\mathsf{sd}}\concat C_{\mathsf{sd}}\concat \mathit{tag},\\
C_{\mathsf{pt}}&\gets
  \AEAD.\Enc(mk,n_{\mathsf{pt}},m^\star,\mathsf{aad}_{\mathsf{pt}}),\\
\sigma&\gets \mathrm{Sign}_{I_i}(S_i.\gid,S_i.\epoch,n_{\mathsf{sd}},
  n_{\mathsf{pt}},C_{\mathsf{sd}},C_{\mathsf{pt}},\mathit{tag}),\\
C&\gets (S_i.\gid,S_i.\epoch,n_{\mathsf{sd}},n_{\mathsf{pt}},
  C_{\mathsf{sd}},C_{\mathsf{pt}},\mathit{tag},\sigma).
\end{aligned}
\]
Set $S_i'\gets S_i[ck_i\leftarrow ck_i',
\mathrm{generation}\leftarrow S_i.\mathrm{generation}+1]$ and persist $S_i'$
before releasing $C$.

\textbf{3. Receiver authenticates the envelope and opens sender data.}
Reject unless $C$ matches the receiver's group identifier and group epoch and
the sender signature verifies for the current roster. Then compute
\[
sd'\gets
  \AEAD.\Dec(S_j.k^{\mathsf{meta}},n_{\mathsf{sd}},C_{\mathsf{sd}},
    S_j.\gid\concat S_j.\epoch).
\]
Parse $sd'$ as $(i',g',\mathit{reuseGuard}')$ and set
$\mathsf{rk}_{\mathsf{replay}}'\gets(i',g')$. Reject if sender-data decryption
fails, if the sender leaf or generation is malformed for the current roster, or
if $\mathsf{rk}_{\mathsf{replay}}'$ is already present in the replay window.

\textbf{4. Receiver resolves the sender-chain key and decrypts the payload.}
\[
\begin{aligned}
(mk',S_j^\star)&\gets \mathrm{ResolveSenderKey}(S_j,sd'),\\
m^\star{}'&\gets
  \AEAD.\Dec(mk',n_{\mathsf{pt}},C_{\mathsf{pt}},
    S_j.\gid\concat S_j.\epoch\concat C_{\mathsf{sd}}\concat \mathit{tag}).
\end{aligned}
\]
Reject if key resolution or payload decryption fails.

\textbf{5. Receiver applies policy and commits receive state.}
Open the payload semantics to obtain $m'$, reject unless
$\mathrm{PolicyAccepts}(S_j.\policy,m^\star{}',\mathit{tag})$ holds. Set
$S_j'\gets \mathrm{CommitReceiveState}(S_j^\star,sd')$, which marks
$\mathsf{rk}_{\mathsf{replay}}'$ as seen, and persist $S_j'$ before exposing
$m'$.
Return $(S_i',C,S_j',m')$.
\end{algorithm}
\endgroup
\FloatBarrier

\section{Message Layer}

TreeKEM Commit messages are comparatively expensive and are used for group-epoch changes.
Inside one group epoch, application messages use sender chains. At group-epoch installation
the base sender chain for leaf $i$ is reset from the accepted group-epoch secret;
for sender leaf $i$ and generation $g$:
\begin{align}
  ck^e_{i,0} &=
  \HKDF\text{-}\mathrm{Expand}\bigl(\epochsec_e,32,\notag\\
  &\qquad \domain{Aura-Group-SenderChain}\concat \gch_e
    \concat \mathrm{le}_{32}(i)\bigr),
    \label{eq:sender-chain-reset}\\
  mk_{i,g} &= \HKDF\text{-}\mathrm{Expand}
    (ck_{i,g},32,\domain{Aura-Group-Msg}),\\
  \widetilde{ck}_{i,g+1} &= \HKDF\text{-}\mathrm{Expand}
    (ck_{i,g},32,\domain{Aura-Group-Chain}),\\
  ck_{i,g+1} &=
  \begin{cases}
    \widetilde{ck}_{i,g+1}, & \text{default policy},\\
    \mathrm{BLAKE2bMAC}_{\domain{Aura-B2Chain}}
      (\widetilde{ck}_{i,g+1}), & \Shield{}.
  \end{cases}
\end{align}
The payload is padded and encrypted under $mk_{i,g}$. Sender metadata
$(i,g,\mathit{reuseGuard})$ is encrypted under the group-epoch metadata key
$k^{\mathsf{meta}}_e$. The message signature covers the group identifier,
group-epoch number, encrypted sender data, encrypted payload, nonces, and franking tag.

\subsection{Replay Rejection}

Each sender has a replay window with a high-water generation and a bitmap.
A duplicate generation inside the window is rejected. A generation too far
behind the window is also rejected. The skipped-key cache is bounded per
sender and globally; \Shield{} tightens the per-sender bound.

\subsection{Franking, Sealed Content, and Expiration}

For a frankable message the sender samples a franking key $fk$ and publishes:
\begin{equation}
  \mathit{tag} =
  \begin{cases}
    \HMAC(fk,\mathit{content}\concat\mathit{sealedContent}), & \text{frankable},\\
    \varepsilon, & \text{otherwise}.
  \end{cases}
  \label{eq:franking}
\end{equation}
The tag is visible in the application message, while $fk$ is inside the
end-to-end encrypted payload. A recipient who reports a message can disclose
$(fk,\mathit{content},\mathit{sealedContent},\mathit{tag})$ for that message
only. The relay cannot verify or read unreported messages because it does not
know $fk$; this follows the message franking goal of selective accountable
disclosure~\cite{franking}.

\begin{figure}[tbp]
\centering
\resizebox{0.98\textwidth}{!}{%
\begin{tikzpicture}[
  font=\small,
  role/.style={draw,rounded corners=2pt,minimum width=2.55cm,
    minimum height=0.72cm,align=center,fill=gray!7},
  note/.style={draw,rounded corners=2pt,align=center,fill=blue!5,
    font=\scriptsize,text width=2.65cm,minimum height=0.72cm},
  relaynote/.style={note,fill=orange!8},
  reportnote/.style={note,fill=green!7},
  arrow/.style={-Latex,thick},
  weak/.style={-Latex,thick,dashed,gray!75}
]
\node[role] (sender) at (0,0) {Sender};
\node[role] (relay) at (3.65,0) {Relay};
\node[role] (recipient) at (7.3,0) {Recipient};
\node[role] (moderator) at (10.95,0) {Moderator};

\node[note] (seal) at (0,-1.18) {$fk$ sealed\\inside E2E payload};
\node[relaynote] (opaque) at (3.65,-1.18) {stores opaque\\$(C,\mathit{tag})$};
\node[reportnote] (open) at (7.3,-1.18) {can disclose\\one message tuple};
\node[reportnote] (verify) at (10.95,-1.18) {verifies one\\reported message};

\draw[arrow] (sender) -- node[above,font=\scriptsize,fill=white,inner sep=1pt]
  {$C,\mathit{tag}$} (relay);
\draw[arrow] (relay) -- node[above,font=\scriptsize,fill=white,inner sep=1pt]
  {$C,\mathit{tag}$} (recipient);
\draw[weak] (open) -- (verify);
\draw[arrow] (sender) -- (seal);
\draw[arrow] (relay) -- (opaque);
\draw[arrow] (recipient) -- (open);
\draw[arrow] (moderator) -- (verify);
\end{tikzpicture}
}
\caption{Franking disclosure path. The relay stores an opaque commitment, while
a recipient can later reveal exactly one message to a moderator.}
\label{fig:franking-flow}
\end{figure}

Sealed content places an encrypted inner payload behind a visible hint, in the
same relay-blind spirit as sealed-sender metadata minimization~\cite{sealed-sender}.
The seal key is derived from the message key, and the recipient obtains it only
after group decryption. Disappearing messages carry a send timestamp and TTL;
decryption rejects expired messages and timestamps too far in the future.

\begin{figure}[tbp]
\centering
\begin{tikzpicture}[
  font=\small,
  layer/.style={draw,rounded corners=2pt,align=center,minimum width=10.5cm,minimum height=0.78cm},
  arrow/.style={-Latex,thick}
]
\node[layer,fill=gray!8] (wire) {GroupMessage: version, group id, group epoch, application frame};
\node[layer,below=0.45cm of wire,fill=blue!6] (sender) {Encrypted sender data: sender leaf, generation, reuse guard};
\node[layer,below=0.45cm of sender,fill=green!7] (payload) {Encrypted payload: content, policy, franking key, sealed fields, references};
\node[layer,below=0.45cm of payload,fill=orange!8] (sem) {Optional semantics: franking tag, sealed reveal, TTL expiration};
\draw[arrow] (wire) -- (sender);
\draw[arrow] (sender) -- (payload);
\draw[arrow] (payload) -- (sem);
\end{tikzpicture}
\caption{Application-message envelope. Sender data is hidden from passive
observers under the group-epoch metadata key, while the payload uses the sender-chain
message key.}
\label{fig:message-envelope}
\end{figure}

\section{External Join and Relay Model}

External join is supported only when policy permits it. The joining member
uses external key material derived from the current init secret and presents an
authorization signed by an existing member. The authorization binds the group
identifier, group-epoch number, group context hash, external public keys, authorizer leaf,
joiner identity, and expiry. In \Shield{} groups this path is disabled.

The authorization is evaluated against the public state and group context for
the group epoch in which it was issued. A joiner cannot transplant it to a different
tree hash, policy, group identifier, external key pair, or later group epoch without
changing the signed data. Expiry is an authorization-creation boundary: it
prevents a stale authorization from being used to form a new external-join
Commit message. Once a valid external-join Commit message has been formed for the bound group epoch,
offline members may still process that Commit later, because their acceptance
check is the signed Commit message and policy-bound group-epoch transition rather than the
joiner's local wall clock.

The relay model is narrower. The relay may validate envelope structure,
group identifier, group-epoch number, sender signature format, and routing information.
For Commit messages it can maintain tentative delivery state, but it cannot verify the
TreeKEM path secret, confirmation MAC, or final client-side membership state
because it has no group secrets. This separation is important: relay-side
validation is useful for abuse resistance and routing, but it is not a
cryptographic authority.

\begin{figure}[tbp]
\centering
\resizebox{\textwidth}{!}{%
\begin{tikzpicture}[
  font=\small,
  lane/.style={draw,rounded corners=2pt,align=center,minimum width=4.05cm,minimum height=0.82cm},
  relay/.style={lane,fill=orange!9},
  member/.style={lane,fill=blue!6},
  endrelay/.style={lane,fill=orange!15,thick},
  endmember/.style={lane,fill=green!8,thick},
  arrow/.style={-Latex,thick},
  boundary/.style={dashed,gray!70,line width=0.7pt},
  node distance=0.58cm and 0.82cm
]
\node[relay] (r1) at (0,0) {Relay lane: size, version, group id, group epoch};
\node[relay,right=of r1] (r2) {signature syntax\\routing target};
\node[endrelay,right=of r2] (r3) {tentative routing\\state only};

\node[member] (m1) at (0,-1.8) {Member lane: decrypt sender data};
\node[member,right=of m1] (m2) {process UpdatePath\\tree hash, policy, MAC};
\node[endmember,right=of m2] (m3) {accepted group epoch\\or accepted message};

\draw[arrow] (r1) -- (r2);
\draw[arrow] (r2) -- (r3);
\draw[arrow] (m1) -- (m2);
\draw[arrow] (m2) -- (m3);
\draw[boundary] ($(r1.south west)+(0,-0.36)$) -- ($(r3.south east)+(0,-0.36)$);
\end{tikzpicture}%
}
\caption{Relay/member authority split. Relay checks are structural and useful for routing, while cryptographic acceptance remains a member-side state transition.}\label{fig:relay-authority}
\end{figure}

\begin{table}[tbp]
\centering
\caption{Relay-visible checks versus member-side checks.}\label{tab:relay}
\small
\begin{tabularx}{\textwidth}{@{}p{0.24\textwidth}YY@{}}
\toprule
Check & Relay side & Member side \\
\midrule
Envelope format & Version, size, group identifier, group-epoch number, payload type. & Same, plus decryption under local group-epoch keys. \\
Commit structure & Group-epoch increment, group identifier, syntactic proposal shape, signature binding. & Proposal validity, TreeKEM path processing, tree hash, confirmation MAC. \\
Application message & Sender identity binding from authenticated context, signature syntax, routing target. & Sender-data decrypt, replay window, sender-chain key, AEAD payload decrypt, policy checks. \\
Membership state & Tentative routing roster. & Authoritative ratchet-tree state and accepted group epoch. \\
\bottomrule
\end{tabularx}
\end{table}

\section{Security Argument}\label{sec:security}

The security argument is modular rather than fully composable. It separates
the cryptographic assumptions used by individual mechanisms from the protocol
invariants that connect those mechanisms inside one accepted group epoch. This is a
defined claim boundary: the argument does not prove arbitrary asynchronous
group schedules. It argues that, for a locally accepted transition, the tree
update, policy context, confirmation MAC, sender-chain state, replay window,
and message semantics fit together in the intended way.

\begin{table}[tbp]
\centering
\caption{Assumption use in the informal security argument.}\label{tab:assumptions}
\small
\begin{tabularx}{\textwidth}{@{}p{0.25\textwidth}YY@{}}
\toprule
Mechanism & Assumption or invariant & Used for \\
\midrule
X25519 path component & Classical DH hardness before a CRQC and fresh post-compromise keys. & Hybrid path secrecy and local recovery when the classical component survives. \\
ML-KEM path component & IND-CCA2 security of ML-KEM-768. & Path secrecy against harvest-now-decrypt-later when the KEM component survives. \\
HKDF schedule & Extractor/PRF behavior with domain-separated labels. & Combining previous init, commit secret, group-epoch context, and sender-chain derivation. \\
AES-256-GCM-SIV & Authenticated encryption with misuse resistance under bounded nonce misuse. & Payload, sender-data, Welcome, and sealed-state confidentiality/integrity. \\
Ed25519 signatures & Unforgeability of identity, Commit-message, and message signatures. & Key-package authenticity, Commit origin, and sender binding. \\
Local erasure & Old path and group-epoch secrets are erased; any retained delayed-message keys are bounded and tagged by group epoch. & Forward secrecy, delayed delivery, and post-compromise recovery boundaries. \\
State continuity & Commit/message acceptance is atomic with durable replay and sealed-state updates. & Preventing rollback from turning old ciphertexts into fresh events. \\
\bottomrule
\end{tabularx}
\end{table}

The central invariant is that there is a single accepted context for a
member at a group epoch: the tuple $(\gid,e,\treehash_e,\policy)$ determines
$\gch_e$, and $\gch_e$ determines the group-epoch keys and confirmation MAC. If any
component changes, honest clients either derive different keys or fail
confirmation. The second invariant is membership-indexed key availability:
only leaves represented in the accepted ratchet tree, and not removed by the
accepted Commit, obtain the path secret needed for the next group epoch. The third
invariant is sender-indexed message state: once a group epoch is accepted, each
sender has an independent chain and replay window, so duplicate generations do
not become fresh messages. The fourth invariant is state continuity: rejected
Commit messages and rejected ciphertexts do not partially advance group-epoch counters,
sender chains, replay windows, or sealed-state freshness markers.

The positive claims below should therefore be read as one-step accepted-state
claims. They are not global liveness claims about delivery, deniability claims,
or a universal-composability theorem for arbitrary group schedules. Within that
boundary, the argument identifies the exact state transition at which fresh
hybrid path material, policy consensus, and sender-chain reset enter the
protocol state.

\begin{claim}[Hybrid group-epoch recovery]
Assume honest erasure of old path secrets and group-epoch secrets. If an accepted
Commit after compromise contains fresh \DH{} and \KEM{} material generated by
an uncompromised member, then the new group-epoch secret is pseudorandom provided at
least one of the hybrid wrapping components remains secure and HKDF behaves as
a randomness extractor and PRF.
\end{claim}

\noindent
The argument follows the hybrid-combiner principle. Path secrets are delivered
under a wrapping key derived from both a \DH{} shared value and a \KEM{} shared
secret. Replacing one surviving component by random changes the wrapping key
to a pseudorandom key. The commit secret then feeds HKDF-Extract with the
previous init secret; the group-epoch context binds the accepted tree and policy.
If the previous init secret was compromised, recovery depends on the fresh
commit secret. If the update path was generated before compromise or by a
compromised member, this claim does not apply. Thus the recovery
point is the first accepted honest group-epoch update whose path material was
sampled after the compromise event.

\begin{claim}[Policy consensus]
Two members that accept the same Commit but use different serialized security
policies derive different group context hashes except with collision
probability for \SHA{}. Consequently, they derive different group-epoch keys and fail
confirmation or later AEAD checks.
\end{claim}

\noindent
This claim is the core reason to treat policy as cryptographic input. A relay
or application cannot silently downgrade \Shield{} to a weaker policy while
preserving the same group-epoch transcript.
The same reasoning covers sender-key limits, mandatory franking, and external
join policy because those fields are serialized into the policy bytes before
hashing. A member that accepts the weaker policy enters a different group-epoch
state from a member that accepts the stronger one; the mismatch is detected by
confirmation or by later message decryption failure.

\begin{claim}[Removed-member exclusion]
Let member $R$ be removed by an accepted Commit from group epoch $e$ to
$e+1$. If $R$ does not receive a valid path secret for group epoch $e+1$ and
the remaining members erase old group-epoch keys and message keys as specified,
then $R$ cannot derive sender-chain message keys for group epoch $e+1$ from
its old group-epoch-$e$ state.
\end{claim}

\noindent
The public Commit may reveal the proposal that removes $R$, but the new group-epoch
secret is derived from the previous init secret and the fresh commit secret
recovered through the update path. The update path is addressed to the
co-path resolution of current recipients, not to removed leaves. Sender-chain
bases are reset from the new group-epoch secret. Consequently, old sender-chain
state does not extend across the removal boundary. This is a membership-interval
exclusion statement, not a claim that already-delivered group-epoch-$e$ ciphertexts
become undecryptable.

If $R$ is later added again, the admission is a new group-epoch transition with a new
key package or external-join authorization, a new leaf position as accepted by
the tree, and fresh sender-chain bases derived from the re-addition group epoch. The
new membership interval does not bridge the removal gap. In particular,
successful re-addition gives $R$ access to group epochs at and after the new Welcome
or accepted external join, not to application messages from group epochs in which
$R$ was absent.

\begin{claim}[Replay rejection]
Within one group epoch, a valid application ciphertext with the same sender leaf and
generation is accepted at most once by a receiver that preserves its replay
window state.
\end{claim}

\noindent
The receiver decrypts sender data under the group-epoch metadata key, verifies the
sender signature, checks the replay window, obtains the sender-chain message
key, and only then marks the generation as seen. Modified replays fail at
signature, AEAD, franking, or replay checks depending on the mutation.
Persistence matters: if a receiver restores old sealed state without an
external anti-rollback floor, replay windows may also roll back. The security
claim therefore requires sealed-state anti-rollback metadata to be treated as
deployment state rather than as a cosmetic serialization field.

\begin{claim}[Delayed-message state separation]
If delayed application messages from previous group epochs are supported, the retained
compatibility state is bounded, tagged by group epoch, and message-key scoped. Accepting
such a delayed message cannot roll the current group epoch backward, reinstall
an old tree, or revive a removed sender-chain state.
\end{claim}

\noindent
This separates asynchronous delivery from rollback. A member may keep a small
cache of skipped message keys or metadata-opening material for messages that
were already authorized in their source group epoch. That cache is not the previous
group-epoch secret, not the previous init secret, and not authority to process new
Commit messages under an old tree. The message is checked against the group identifier,
source group epoch, sender leaf, generation, replay window, and policy semantics for
the group epoch that produced the key. If the cache is absent or exhausted, the
message is rejected; the implementation does not synthesize old state from the
current group epoch.

\begin{claim}[Franking accountability]
If a moderator verifies a disclosed tuple
$(fk,\mathit{content},\mathit{sealedContent},\mathit{tag})$ for an accepted
frankable message, then the tag corresponds to the franking key and content
fields carried in an end-to-end encrypted payload accepted by a group member,
except with the forgery probability of the MAC.
\end{claim}

\noindent
The relay sees the franking tag before any report, but it does not know the
franking key because $fk$ is carried inside the encrypted payload. A recipient
can disclose $fk$, the relevant content, the sealed-content field, and the tag
for one message. The moderator then checks the MAC relation for that message
only. This does not give the moderator keys for unrelated messages, and it does
not let the relay verify ordinary unreported traffic.

\subsection{Failure Cases Outside the Claim Boundary}

Several important failure cases are outside the positive claims.
If all live members that can update the tree are compromised and later no
honest member performs an update, there is no post-compromise recovery event.
If clients accept unsigned or incorrectly bound key packages, then group
membership authentication has already failed before TreeKEM processing begins.
If a deployment permits policy changes without treating them as explicit
Commit messages, then the policy-consensus argument no longer applies. If
sealed state can be rolled back by the local platform, replay windows and
skipped-key budgets can also be rolled back. These conditions are stated as
part of the interface contract for the protocol state machine.

\section{Scoped Symbolic Model}

The construction is accompanied by the scoped \ProVerif{} model
\texttt{formal/proverif/aura\_group.pv}~\cite{blanchet2016}. It can be run
with \texttt{make -C formal proverif-group}. The model covers one accepted
hybrid update path and one application message inside the resulting group
epoch. The network adversary controls the public channel, can replay and
substitute public messages, and sees the hybrid path wrapper and application
ciphertext; honest acceptance is represented by confirmation-MAC and
correspondence events. The model abstracts away full asynchronous group
scheduling, proposal queues, delivery-service fairness, multiple concurrent
senders, and multi-device state.

\begin{table}[tbp]
\centering
\caption{Verification scope and non-claims for \Aura{} Group.}
\label{tab:group-verification-scope}
\footnotesize
\setlength{\tabcolsep}{2pt}
\begin{tabularx}{\textwidth}{@{}L{0.22\textwidth}L{0.36\textwidth}Y@{}}
\toprule
Evidence source & In scope & Explicitly out of scope \\
\midrule
Protocol argument & Accepted epoch secrecy, hybrid recovery boundary, policy consensus, removed-member exclusion, replay rejection, delayed-message separation, and franking accountability. & Universal composability for all asynchronous group schedules, delivery fairness, and legal or product moderation semantics. \\
\ProVerif{} model & One accepted hybrid update path, one resulting group epoch, one application message, and six secrecy/correspondence queries. & Full ratchet-tree traversal, all co-path recipients, proposal queues, skipped-key caches, replay bitmaps, and multiple concurrent senders. \\
Reference implementation & Parsing, rejection paths, transactional state movement, persistence, replay windows, external-join authorization, \Shield{} policy enforcement, franking, sealed state, and benchmarked cost centers. & A substitute for reduction-style proof, a proof of resource-exhaustion resistance, or evidence for unimplemented deployment policies. \\
Benchmarks & Cost characterization for group creation, admission, application messages, group updates, and sealed-state import/export. & Security assumptions, production capacity guarantees, or mobile-device performance claims. \\
\bottomrule
\end{tabularx}
\end{table}

\begin{table}[tbp]
\centering
\caption{\ProVerif{} group model: scoped queries discharged.}\label{tab:group-proverif}
\small
\begin{tabularx}{\textwidth}{@{}p{0.16\textwidth}p{0.28\textwidth}p{0.14\textwidth}Y@{}}
\toprule
Query & Model event or secrecy target & Result & Meaning \\
\midrule
Q1 & \texttt{path\_secret} secrecy & true & The fresh TreeKEM path secret is not learned from the hybrid path wrapper. \\
Q2 & \texttt{prev\_init\_secret} secrecy & true & The previous init-secret contribution remains hidden in the accepted transition. \\
Q3 & \texttt{app\_secret} secrecy & true & The challenge payload within the accepted group epoch is not learned by the network adversary. \\
Q4 & \texttt{CommitAccepted} $\Rightarrow$ \texttt{CommitCreated} & true & An accepted Commit corresponds to an honest policy-bound Commit event. \\
Q5 & \texttt{MessageAccepted} $\Rightarrow$ \texttt{MessageSent} & true & An accepted challenge message corresponds to an honest send under the same group-epoch key. \\
Q6 & \texttt{FrankingVerified} $\Rightarrow$ \texttt{FrankingIssued} & true & A verified franking disclosure corresponds to an issued franking tag. \\
\bottomrule
\end{tabularx}
\end{table}

The six queries check the two composition points that the construction relies on:
the Commit-to-epoch transition, where a path secret becomes a group-epoch
secret under a policy-bound context, and the epoch-to-message transition,
where the accepted epoch key becomes a message and franking key. The model
represents the hybrid path wrapper, the policy-bound confirmation MAC, the
application payload, and franking correspondence explicitly. Full tree
traversal, all co-path recipients, skipped-key cache behavior, replay bitmaps,
and arbitrary proposal interleavings remain outside this one-epoch model.
This is a scoped positive result: the model tests the intended cryptographic
boundary rather than replacing a future concurrent TreeKEM proof.

\section{Reference Implementation Validation}

The reference-implementation validation evidence is summarized in
\cref{tab:evidence}. It is organized by validation class:
lifecycle transitions, member exclusion, external-join authorization,
message-layer misuse cases, higher-level message semantics, sealed-state
persistence, and cost characterization. These checks show that the executable
state machine enforces the same acceptance and rejection boundaries described
by the specification.

\begin{table}[tbp]
\centering
\caption{Reference-implementation validation evidence for the group construction.}\label{tab:evidence}
\small
\begin{tabularx}{\textwidth}{@{}p{0.26\textwidth}YY@{}}
\toprule
Area & Evidence & Security relevance \\
\midrule
TreeKEM update & Add, remove, update, interleaving, and parent-hash validation cases. & Group-epoch state evolves through authenticated update paths. \\
Removed member exclusion & Rejection of post-removal decrypt attempts and stale captured ciphertexts. & Old members do not obtain post-removal group-epoch keys. \\
External join & Rejection of stale, expired, tampered, wrong-state, or policy-blocked join attempts. & Join surface is explicit and policy controlled. \\
Message layer & Replay, ciphertext-tamper, and sender-authentication rejection cases. & Application messages are bound to group epoch, sender, and generation. \\
Message semantics & Validation of sealed, expiring, frankable, edited, and deleted message semantics. & Higher-level semantics are enforced after cryptographic decrypt. \\
Sealed state & Rejection of malformed, wrong-key, wrong-identity, tampered, and rolled-back state. & Local state persistence has integrity and external freshness requirements. \\
Performance & Measurements of group creation, admission, message processing, updates, and state sealing. & Cost measurements are engineering context, not security assumptions. \\
\bottomrule
\end{tabularx}
\end{table}

The validation suite exercises the construction at several layers: external
interface behavior, lower-level TreeKEM invariants, key-package validation,
key-schedule domain separation, message-layer misuse cases, time-dependent
semantics, persistence checks, and adversarial rejection paths. The evidence is
reported by state-machine coverage: the executable transitions are checked
against the same acceptance and rejection boundaries specified by the
construction.

The validation evidence covers hybrid path-secret processing, parent-hash
validation, path-key rotation, pre-join and post-removal confidentiality
boundaries, re-addition gap protection, per-group-epoch replay rejection,
external-join authorization binding, \Shield{} policy enforcement, franking
tamper rejection, and sealed-state integrity. This evidence supports the
protocol description at the engineering boundary where parsing, persistence,
replay handling, and rejection behavior are observable.

Benchmarks are reported separately from security evidence. They measure
the principal cost centers of the construction: creating a group, admitting a
member, encrypting and decrypting group messages, performing an update, and
exporting or importing sealed group state. Benchmark timings are not used as
assumptions; the measurements are engineering context.

\subsection{Reproducibility Notes}

A reviewer can examine the construction through four views: the protocol
description, the reference implementation, the validation corpus, and the
symbolic model. The manuscript gives the protocol-level equations and claim
boundary; the implementation realizes tree state, TreeKEM wrapping, Commit
messages, key packages, Welcome processing, sender keys, and key schedule
derivation; the validation material and symbolic model exercise these layers
from complementary perspectives. Benchmark tables and validation counts are
included for reproducibility and engineering context.

The artifact workflow records the exact source revision, operating system, CPU,
Rust toolchain, dependency lockfile, dirty-worktree state, validation commands,
formal-model logs, and benchmark summaries. This is the intended
reproducibility boundary: a reviewer can rerun the symbolic model, validation
suite, and group benchmark targets, but the resulting logs remain engineering
evidence rather than additional cryptographic claims.

\section{Comparison}

\begin{table}[tbp]
\centering
\caption{Design positioning against common group-messaging design points.}\label{tab:comparison}
\small
\begin{tabularx}{\textwidth}{@{}p{0.23\textwidth}YYY@{}}
\toprule
Property & MLS 1.0 baseline & Sender-key style groups & \Aura{} Group profile \\
\midrule
Group-epoch structure & Standardized TreeKEM epochs and group context in RFC 9420. & Often pairwise admission plus symmetric sender-key traffic. & MLS-inspired epochs with hybrid TreeKEM update-path wrapping. \\
Post-quantum material & Not specified by the RFC 9420 baseline. & Depends on the system's pairwise setup and update policy. & \DH{}+\KEM{} material carried in membership-changing update paths. \\
Policy as key input & Group context is central; application policy profiles are separate. & Commonly enforced above the cryptographic group state. & Serialized policy is hashed into the epoch context and subkey schedule. \\
Within-epoch messages & Application secrets derive from the accepted epoch state. & Efficient sender chains. & Per-member sender chains with encrypted sender data and epoch binding. \\
Accountability hooks & Not the main protocol goal. & Product-specific. & Franking and sealed-content hooks bound to accepted epoch policy. \\
Relay authority & Delivery service is not a group-secret holder. & Service role varies by system. & Relay performs structural checks; clients decide cryptographic acceptance. \\
\bottomrule
\end{tabularx}
\end{table}

\subsection{Interpretation}

The comparison in \cref{tab:comparison} positions the protocol rather than
ranking systems on a single axis. MLS is the mature standard reference for
TreeKEM-based group messaging. Sender-key style designs are efficient for
application traffic. \Aura{} Group combines those design principles in a
narrower profile: use a tree for epoch-changing membership updates, carry
hybrid \DH{}+\KEM{} material in the update path, and bind sender-chain message
processing to the accepted epoch policy.

The main tradeoff is complexity at the group-epoch boundary. A hybrid update path
must carry more public material and more ciphertext material than a purely
classical update. It also requires careful binding between recipient
resolution, node index, tree hash, and policy context. The benefit is that a
fresh honest group-epoch update can incorporate post-quantum material into the group
state itself, rather than confining the post-quantum contribution to admission.

The policy-binding design also has a deployment cost. If a product wants
mutable policy, such as toggling franking or changing skipped-key limits, that
policy change must be represented as an explicit group-epoch transition. This is
less convenient than a server-side flag, but it yields an explicit
cryptographic downgrade boundary for protocol participants. In high-risk
groups that boundary is preferable: policy is accepted when the group epoch is
accepted, not when a relay says it has changed.

The scientific point is therefore not the invention of TreeKEM, sender keys, or
message franking in isolation. It is the composition boundary for a concrete
post-quantum group profile: a single accepted group epoch determines tree state,
policy, metadata protection, sender-chain bases, replay state, and accountable
message semantics. This is the boundary that the equations, algorithms,
symbolic model, and implementation evidence are organized to inspect.

\section{Limitations}

\begin{table}[H]
\centering
\caption{Ranked limitations and deployment effect.}\label{tab:group-limitations}
\footnotesize
\setlength{\tabcolsep}{2pt}
\begin{tabularx}{\textwidth}{@{}L{0.13\textwidth}L{0.25\textwidth}Y@{}}
\toprule
Severity & Limitation & Effect on the claim boundary \\
\midrule
High & Classical identity and Commit signatures & Identity and Commit-message authentication use Ed25519. Hybrid TreeKEM secrecy is post-quantum only for the key-establishment material; migration to ML-DSA or another post-quantum signature scheme is separate work. \\
High & Concurrent group proof boundary & The \ProVerif{} model covers one accepted group-epoch transition and one application message. Complete group-security claims require explicit tree state, proposal queues, external joins, replay windows, and concurrency at the level expected in formal MLS analyses~\cite{treesync,meier2013}. \\
High & Platform rollback and delivery fairness & The protocol rejects stale or malformed state, but it cannot force relay delivery or prevent a platform from rolling back local sealed state without an external freshness floor. Replay and delayed-message claims therefore depend on durable state continuity. \\
Medium & Single-device model & Multi-device fan-out, per-device membership, device churn, and cross-device state repair require a separate membership and recovery model. \\
Medium & Relay validation is structural & The relay can enforce version, size, routing, and syntactic checks, but it cannot validate TreeKEM secrets or act as authoritative group membership state. \\
Medium & Mutable policy requires explicit Commit semantics & The current construction treats policy as part of accepted group state. Dynamic policy changes require protocol-level Commit semantics and downgrade analysis, not a server-side flag. \\
Medium & No deniability proof & Message signatures and franking support sender binding and accountable disclosure. The manuscript does not claim deniable authentication for group application messages. \\
Medium & No full resource-exhaustion proof & The implementation bounds parsing, skipped-key storage, and replay windows, but adversarial memory/CPU exhaustion under hostile traffic is an engineering analysis outside the present model. \\
\bottomrule
\end{tabularx}
\end{table}

\section{Conclusion}\label{sec:conclusion}

\Aura{} Group specifies a compact hybrid group-messaging protocol in which
membership updates, policy agreement, and application-message processing are
all anchored at the accepted group epoch. The central design choice is to put
post-quantum material in ordinary \TreeKEM{} update paths and to put serialized
security policy in the group context that drives key derivation and
confirmation. As a result, a membership-changing Commit does more than update
the tree: it installs the cryptographic interval in which sender chains,
metadata protection, \Shield{} processing, franking, sealed state, replay
windows, and expiration checks are interpreted.

The paper gives an inspectable design point and a stated evidence boundary for
post-quantum confidential group state. The informal argument covers hybrid epoch
recovery, policy consensus, removed-member exclusion, replay rejection,
delayed-message separation, and franking accountability under explicit erasure
and state-continuity assumptions. The scoped \ProVerif{} model checks six
queries for one accepted update path and one application message, and the
implementation validation exercises the parser, state machine, rejection paths,
external-join controls, \Shield{} policy, franking, sealed persistence, and
benchmark cost centers. The result is a protocol-level account of how
post-quantum path material and application-security policy enter the same
accepted group state, with the remaining concurrency and deployment obligations
made explicit rather than hidden in the implementation.

The remaining work is correspondingly clear: post-quantum
authentication, a concurrent tree proof for arbitrary proposal schedules,
multi-device group state, delivery-service fairness, platform anti-rollback
anchors, mobile-class benchmarks, and hostile-traffic resource analysis. Within
the scope stated here, \Aura{} Group offers a concrete protocol construction for
policy-bound hybrid TreeKEM epochs and efficient group messages under an
auditable post-quantum confidentiality boundary.

\section*{Data and Software Availability}

The source code, validation material, symbolic models, benchmark summaries,
reproducibility scripts, and manuscript sources are maintained in the
accompanying source repository:
\url{https://github.com/oleksandrmelnychenko/aura-protected-protocol-rs}.
The accompanying artifact contains the exact source snapshot used for the
implementation, symbolic-model, validation, and benchmark results, together
with the generated \texttt{versions.txt}, \texttt{MANIFEST.txt},
\texttt{SHA256SUMS}, formal-model logs, benchmark summaries, and validation
outputs. These files are the reproducibility record for the results in this
manuscript; later repository revisions are treated as continuing development
rather than as the cited artifact.

\section*{Funding}

The author declares that this research received no external funding.

\section*{CRediT Author Statement}

Oleksandr Melnychenko: conceptualization, methodology, software,
validation, symbolic modeling, investigation, writing--original draft,
writing--review and editing, and visualization.

\section*{Declaration of Competing Interest}

The author declares that there are no known competing financial
interests or personal relationships that could have appeared to
influence the work reported here.

\section*{Declaration of Generative AI and AI-assisted technologies in the writing process}

OpenAI Codex was used during manuscript preparation to assist with
English-language editing, consistency checks, and minor \LaTeX{}
typesetting support. OpenAI Codex was not used as an authoritative
source for technical claims, references, equations, implementation
results, or experimental measurements. All protocol descriptions,
equations, references, tables, implementation details, and experimental
results were independently reviewed, verified, and approved by the
author, who takes full responsibility for the content of the article.

\bibliographystyle{elsarticle-num}
\begin{thebibliography}{99}

\bibitem{melnychenko2026aura}
O.~Melnychenko.
\newblock Aura Protocol: Hybrid Post-Quantum Confidentiality for Secure
Messaging via Per-Ratchet KEM Rotation.
\newblock Companion manuscript, 2026.
\newblock Manuscript under preparation.

\bibitem{melnychenko2026opaque}
O.~Melnychenko.
\newblock Hybrid Post-Quantum Authentication from OPAQUE and ML-KEM-768:
Construction, Formal Analysis, and a Reproducible Implementation.
\newblock Companion manuscript, 2026.
\newblock Manuscript under preparation.

\bibitem{melnychenko2026opaqueformal}
O.~Melnychenko.
\newblock Formal Verification Report---Hybrid PQ-OPAQUE Protocol.
\newblock Companion formal-evidence report, 2026.
\newblock Technical report under preparation.

\bibitem{jarecki2018}
S.~Jarecki, H.~Krawczyk, and J.~Xu.
\newblock {OPAQUE}: An asymmetric {PAKE} protocol secure against
pre-computation attacks.
\newblock In \emph{EUROCRYPT 2018}, pages 456--486. Springer, 2018.
\newblock \url{https://doi.org/10.1007/978-3-319-78372-7_15}

\bibitem{ietf-opaque}
D.~Bourdrez, H.~Krawczyk, K.~Lewi, and C.~A.~Wood.
\newblock The {OPAQUE} Augmented Password-Authenticated Key Exchange
({aPAKE}) Protocol.
\newblock RFC 9807, 2025.
\newblock \url{https://www.rfc-editor.org/info/rfc9807}

\bibitem{svystun2025dytam}
S.~Svystun, L.~Scis{\l}o, M.~Pawlik, O.~Melnychenko, P.~Radiuk,
O.~Savenko, and A.~Sachenko.
\newblock DyTAM: Accelerating wind turbine inspections with dynamic UAV
trajectory adaptation.
\newblock \emph{Energies}, 18(7):1823, 2025.
\newblock \url{https://api.crossref.org/works/10.3390/en18071823}

\bibitem{svystun2024thermal}
S.~Svystun, O.~Melnychenko, P.~Radiuk, O.~Savenko, A.~Sachenko, and
A.~Lysyi.
\newblock Thermal and RGB images work better together in wind turbine damage
detection.
\newblock \emph{International Journal of Computing}, 23(4):526--535, 2024.
\newblock \url{https://doi.org/10.47839/ijc.23.4.3752}

\bibitem{lysyi2025fire}
A.~Lysyi, A.~Sachenko, P.~Radiuk, M.~Lysyi, O.~Melnychenko, O.~Ishchuk,
and O.~Savenko.
\newblock Enhanced fire hazard detection in solar power plants: an integrated
UAV, AI, and SCADA-based approach.
\newblock \emph{Radioelectronic and Computer Systems}, (2):99--117, 2025.
\newblock \url{https://doi.org/10.32620/reks.2025.2.06}

\bibitem{rfc9420}
R.~Barnes, B.~Beurdouche, R.~Robert, J.~Millican, E.~Omara, and K.~Cohn-Gordon.
\newblock The Messaging Layer Security (MLS) Protocol.
\newblock RFC 9420, 2023.
\newblock \url{https://www.rfc-editor.org/rfc/rfc9420}

\bibitem{treekem}
K.~Bhargavan, R.~Barnes, and E.~Rescorla.
\newblock TreeKEM: Asynchronous Decentralized Key Management for Large Dynamic
Groups.
\newblock Protocol proposal for Messaging Layer Security, 2018.
\newblock \url{https://mailarchive.ietf.org/arch/msg/mls/v1CY0jFAOVOHokB4DtNqS__tX1o/4/}

\bibitem{treesync}
T.~Wallez, J.~Protzenko, B.~Beurdouche, and K.~Bhargavan.
\newblock TreeSync: Authenticated Group Management for Messaging Layer Security.
\newblock USENIX Security Symposium, 2023.
\newblock \url{https://www.usenix.org/conference/usenixsecurity23/presentation/wallez}

\bibitem{nist-pqc}
National Institute of Standards and Technology.
\newblock Post-Quantum Cryptography Standardization.
\newblock \url{https://csrc.nist.gov/projects/post-quantum-cryptography}

\bibitem{fips203}
National Institute of Standards and Technology.
\newblock Module-Lattice-Based Key-Encapsulation Mechanism Standard (ML-KEM).
\newblock FIPS 203, 2024.
\newblock \url{https://doi.org/10.6028/NIST.FIPS.203}

\bibitem{rfc7748}
A.~Langley, M.~Hamburg, and S.~Turner.
\newblock Elliptic Curves for Security.
\newblock RFC 7748, 2016.
\newblock \url{https://www.rfc-editor.org/rfc/rfc7748}

\bibitem{rfc5869}
H.~Krawczyk and P.~Eronen.
\newblock HMAC-based Extract-and-Expand Key Derivation Function (HKDF).
\newblock RFC 5869, 2010.
\newblock \url{https://www.rfc-editor.org/rfc/rfc5869}

\bibitem{rfc8452}
A.~Gueron, A.~Langley, and Y.~Lindell.
\newblock AES-GCM-SIV: Nonce Misuse-Resistant Authenticated Encryption.
\newblock RFC 8452, 2019.
\newblock \url{https://www.rfc-editor.org/rfc/rfc8452}

\bibitem{franking}
P.~Grubbs, J.~Lu, and T.~Ristenpart.
\newblock Message franking via committing authenticated encryption.
\newblock \emph{Advances in Cryptology -- CRYPTO 2017}.
\newblock \url{https://eprint.iacr.org/2017/664}

\bibitem{sealed-sender}
Signal.
\newblock Technology preview: Sealed Sender for Signal.
\newblock Signal Blog, 2018.
\newblock \url{https://signal.org/blog/sealed-sender/}

\bibitem{blanchet2016}
B.~Blanchet.
\newblock Modeling and verifying security protocols with the applied pi calculus
and {ProVerif}.
\newblock \emph{Foundations and Trends in Privacy and Security},
1(1--2):1--135, 2016.
\newblock \url{https://bblanche.gitlabpages.inria.fr/publications/BlanchetFnTPS16.html}

\bibitem{meier2013}
S.~Meier, B.~Schmidt, C.~Cremers, and D.~Basin.
\newblock The {TAMARIN} prover for the symbolic analysis of security protocols.
\newblock In \emph{CAV 2013}, pages 696--701. Springer, 2013.
\newblock \url{https://doi.org/10.1007/978-3-642-39799-8_48}

\end{thebibliography}

\end{document}
